% SOURCE docs/release/2026-10-29/source/docs/writing-repair/2026-10-29/src/3b3c4865206192c5.tex CHAR 0
\chapter{Commercial identification, counterfactuals and the information a bank must supply}
\label{exec13:commercial-identification}

The preceding chapter made acquisition conditional on preserving the scientific
question. This chapter develops that condition through identification,
counterfactuals, heterogeneity, instruments, event histories and policy value.
It then compares the engine and application boundaries that carry those objects.
The method studies establish the targets and counterexamples; the later
product studies use them to choose a particular interface, entitlement and
qualification path. These are related decisions, not independent demonstrations
of the same claim.

A lending manager asks whether offering a temporary rate reduction will improve
twelve-month repayment among existing borrowers. A second manager asks whether
one particular borrower would have repaid without the offer. A third asks how
much of the effect operates through an intermediate change in payment behavior.
These requests share business language and variables, but they demand different
causal information. Purchasing a package with a causal menu does not make them
the same statistical problem. Nor does implementing a common API remove their
different assumptions.

The proposed product should first identify the requested quantity and then
compare the information and computation required to answer it. Commercial
software provides substantive precedents: graphical intervention workflows,
treatment-effect estimators, mediation models and, in at least one inspected
guide, an explicit abduction workflow. Those are capabilities against which the
proposal must be assessed. Their documentation also illustrates why a numerical
answer needs a stated model and why an unsupported question must remain visible.

This comparison draws on a declared set of sixteen commercial families and
27 bounded public document requests, together with the preserved commercial
corpus. It inspects technical documents for BayesiaLab, Stata and Bayes Server;
an Mplus resource index supplies additional leads. Other families retain
specific unresolved comparisons. The literature cutoff remains 7 October 2026.
New document copies were acquired on 13 October, and mutable pages do not
authenticate their exact contents at the cutoff. Stata's inspected PDFs carry
2025 copyright and release-19 data references; the Bayes Server guide names
specific version thresholds. These identities are more informative than an
undated search result, but they do not eliminate the historical-copy limitation.
No licensed vendor engine was executed for this comparison
\citep{exec13_ci_scope}.

\section{Four questions a commercial tool may answer}

Let $A\in\{0,1\}$ denote the offer and $Y(a)$ repayment by a specified date
under offer state $a$. The outcome is defined for every existing loan, including
loans receiving no offer. The first target is
\[
 \tau=\mathbb E\{Y(1)-Y(0)\}.
\]
The population, date, treatment version and accounting meaning are part of
this definition. A target among recipients, conditional on pre-offer history,
or under a sequence of actions needs a different expression. An analysis of
only newly approved loans would additionally face selection and possibly a
different outcome definition; it cannot silently reuse this target.

\subsection{Adjustment and general identification}

An adjustment method starts from a justified observed set $X$. Under consistency,
conditional exchangeability and suitable positivity, it evaluates
\[
 \mathbb E\{Y(a)\}
 =\int \mathbb E[Y\mid A=a,X=x]\,dP(x).
\]
The integral is an identification statement; fitting its regression and
averaging predictions is a subsequent estimation problem. A tool can implement
the latter while relying entirely on the analyst for the former.

A broader identification request supplies a graph class, observed variables,
available observational or experimental distributions and an intervention query.
For example, it may ask for $P(Y\mid\operatorname{do}(A=a))$ in a graph with
unobserved common causes, or for a conditional intervention distribution
$P(Y\mid\operatorname{do}(A=a),Z=z)$. Claims about ID or IDC must name the
algorithm and its supported graph and query classes. Support for one valid
adjustment set does not establish those broader capabilities.

The response contract should distinguish an identified functional, a justified
nonidentification result within the declared model class, and an unknown result
because the available method does not cover the request. Failure of a
back-door search belongs to the third category unless a separate argument
establishes the second. The exact mediator example below shows the practical
importance of this distinction.

\subsection{Individual counterfactuals and mediated effects}

For a borrower who received the offer and repaid, a different query is
\[
 q=\Pr\{Y(0)=0\mid A=1,Y=1\}.
\]
It concerns a relationship between the factual and alternative worlds for
the same unit. Knowing both marginal distributions $P\{Y(0)\}$ and
$P\{Y(1)\}$ need not identify their joint distribution. A structural model
with specified exogenous variables can supply that relationship, but the
relationship is then additional model information. A workflow that infers
those exogenous variables from evidence, changes an action and predicts the
result must preserve the assumptions connecting the worlds.

Mediation makes another distinction. Let $M(a)$ be a mediator, measured after
the offer, and $Y(a,m)$ the outcome under joint interventions. One decomposition
of the total effect uses
\begin{align*}
 \mathrm{NDE}&=\mathbb E[Y(1,M(0))-Y(0,M(0))],\\
 \mathrm{NIE}&=\mathbb E[Y(1,M(1))-Y(1,M(0))].
\end{align*}
Their sum is the total effect under the corresponding composition assumptions.
These are population averages of nested potential outcomes. They are neither
the same target as an individual borrower's $q$ nor automatically identified
by randomizing $A$. The mediator--outcome relation and treatment-induced
confounding still matter. Stata's one-mediator manual makes these assumptions
explicit rather than presenting the word ``counterfactual'' as an unrestricted
capability \citep{exec13_ci_stata_mediate}.

\subsection{Partial identification and uncertainty}

When the information identifies a set rather than a point, the product should
return that set and explain which restrictions would narrow it. This differs
from a confidence interval for a point parameter. A finite-sample confidence
region for an identified set requires an additional sampling argument. A
software comparison must therefore record separately whether a package computes
a point functional, logical bounds, a sensitivity region or statistical
uncertainty, and how failed or empty requests are represented.

\section{BayesiaLab: a documented graphical adjustment workflow}

The inspected graphical-identification chapter explains causal and noncausal
paths, adjustment in a supplied DAG and a particular hidden-common-cause
obstruction. Its discussion supports a comparison of adjustment semantics.
It does not specify a complete general ID/IDC implementation
\citep{exec13_ci_bayesia_criteria}. In particular, its example in which a new
hidden cause defeats adjustment should not be generalized to every graph with
hidden confounding.

The graph-surgery walkthrough is operationally more specific. An analyst draws
the causal graph, associates data columns and state values with its nodes,
estimates the probability tables, and inspects the paths to a selected target.
The intervention menu removes the treatment's incoming arcs for the causal
query. The displayed probabilities then change without treating the action as
ordinary observed evidence. This is a documented supplied-model intervention
workflow, including a visual explanation of the observational/interventional
distinction \citep{exec13_ci_bayesia_surgery}.

For the proposed bank tool, data-to-node association is a substantive capability.
An offer flag recorded after acceptance may not represent the assigned offer;
a repayment field may be available after the declared decision date. A graph
display that preserves node names but loses this mapping can give a convincing
diagram of the wrong question. The proposal should make each node's unit,
clock, state coding and source field reviewable, and make the intervention
status visible alongside the graph. These are concrete requirements suggested
by the documented commercial workflow, not evidence that the local prototype
already implements a comparable visual environment.

The likelihood-matching walkthrough documents another route for the same
example. It uses probabilistic evidence and a \emph{Fix Probabilities} operation
to hold a confounder's marginal distribution fixed while varying treatment
evidence. Its reported example matches the graph-surgery calculation. The
important mathematical question is whether the variables and distributions
being held fixed supply a valid adjustment for the requested target
\citep{exec13_ci_bayesia_matching}. A numerical balancing operation does not
establish that an unmeasured common cause has been removed.

Bayesia's contribution-analysis page uses generated data, a learned network,
expert assumptions about confounder selection and Likelihood Matching to compare
factual and intervention conditions. That supports a documented contribution
workflow. The inspected text does not fully specify a general individual-level
abduction model or the coupling of alternative outcomes. Its linked slides
and recording were not inspected in this comparison
\citep{exec13_ci_bayesia_counterfactual}. The proposed product should therefore
ask whether a request is a population intervention contrast, a contribution
allocation or a conditional individual counterfactual before choosing an adapter.

The reuse recommendation is conditional: treat these workflows as commercial
benchmarks and possible licensed integration targets. Obtain the exact API,
edition, export semantics and deployment rights before promising an adapter.
Porting proprietary implementation is not an available reuse strategy on the
basis of public documentation. Independently implementing a justified
mathematical method requires its own provenance and validation.

\section{Stata: treatment effects and mediation with explicit assumptions}

The inspected \texttt{teffects} introduction documents average treatment effects,
effects among treated subjects and potential-outcome means. It covers regression
adjustment, inverse-probability weighting, augmented weighting, weighted
regression adjustment and matching. Its caveats discuss independent sampling,
conditional independence and overlap, including why unmeasured determinants of
treatment and outcome can invalidate an estimate
\citep{exec13_ci_stata_causal_intro}.

This is a substantial estimator baseline. The proposal needs a defensible reason
to build another estimator rather than a thin, licensed workflow around an
existing one. A possible reason is preservation of the business target,
information clock, graph ambiguity and source provenance across several tools.
That reason becomes persuasive only when a minimal prototype demonstrates the
behavior and a comparison measures the total effort of obtaining an acceptable
answer. Merely presenting additional estimators or claiming a common CLI does
not establish superiority over an established analysis workflow.

For the lending request, an analyst could supply pre-offer covariates and fit
one of the documented treatment-effect estimators. The output then inherits
the asserted covariate sufficiency and the relevant sampling law. The selected
manual's IID discussion does not establish validity for every repeated monthly
loan record or shared-calendar shock. Likewise, an overlap diagnostic concerns
observed treatment support; it cannot certify that an omitted repayment driver
has been measured. The proposed wrapper must record these distinct obligations
instead of converting a successful command exit into an identification pass.

The \texttt{mediate} manual supplies a separate, developed capability. It
describes continuous, binary and count outcome/mediator models, several treatment
types, interactions, direct and indirect effects, potential-outcome means and
variance options. It explicitly addresses unmeasured confounding in the
treatment--outcome, treatment--mediator and mediator--outcome relations, and
excludes treatment-induced mediator--outcome confounding under its stated
identification conditions \citep{exec13_ci_stata_mediate}.

Consider an offer that changes a borrower's monthly payment burden, which then
affects repayment. The offer might also trigger a collections contact that
affects both payment behavior and repayment. That contact is a possible
treatment-induced mediator--outcome confounder. Adding it mechanically to every
regression is not a justification of the natural-effect estimand. The business
request may need a different mediated target or a richer design. A common API
should carry the precise nested expression and its assumptions rather than
label every coefficient product an indirect causal effect.

The inspected manuals support these estimator and mediation comparisons.
They do not establish generic graphical ID/IDC, arbitrary individual
abduction or complete latent-graph coverage. Other Stata commands may address
other endogeneity models; those require their own technical comparisons. The
current product proposal should prioritize matched documented targets and keep
unreviewed extensions explicit. Exact licensing, embedded execution and price
comparisons remain outstanding; no commercial runtime result is reported here.

\section{Bayes Server: intervention, abduction and the need for an exact API contract}

The Bayes Server causal guide distinguishes ordinary evidence from intervention,
describes discrete and continuous variables, and gives a concrete individual
counterfactual workflow. Its example enters factual evidence, uses an
\emph{Abduction} dialog to infer characteristic variables, applies a changed
intervention, and reads the resulting outcome. The guide describes helper
functions from version 10 onward and says the example calculations can be
performed from version 9.4 \citep{exec13_ci_bayesserver_causality}.

This is stronger evidence of a documented abduction workflow than a generic
``what if'' feature label. It reveals another capability gap for the proposal:
the local reference service does not yet offer this general structural-model
workflow. The architecture should represent the exogenous-variable model,
conditioning evidence, action and resulting distribution as separate objects.
It should also retain whether abduction produced a unique value, a distribution
or an unresolved set. A model-specified answer can be useful while remaining
conditional on the model's cross-world assumptions.

The guide also lists back-door and front-door adjustment, instruments,
disjunctive-cause reasoning and do-calculus, stating that many approaches are
available from version 10 onward. That phrasing does not identify which exact
algorithms are implemented, the graph classes for which they are complete, or
the supported failure certificates. Section~\ref{exec14:commercial-contracts} develops the retained Advanced
Causal Query and Abduction guides and API signatures. Their criterion-specific
information resolves this narrower documentation question; complete ID/IDC
and arbitrary counterfactual identification remain unestablished.

An adapter proposal needs more than the guide's terminology. It needs exact
node and edge encodings, latent-variable treatment, allowed interventions,
conditioning behavior, numerical types, export format, method/version identity
and a licensed execution boundary. Until those are inspected and tested, the
recommendation is a focused commercial benchmark candidate. It is not a promise
of immediately reusable functionality or a claim that a general causal
identification engine is absent from the product.

\section{The remaining commercial landscape and explicit omissions}

The Mplus mediation index names counterfactual causal effects, examples and
\texttt{MODEL INDIRECT} material associated with version 7.2. It is a useful
route to original technical sources \citep{exec13_ci_mplus_mediation}. An index
does not supply the assumptions or proofs of the linked papers. Those originals
must be acquired as genuine PDFs and inspected before they support additional
scholarly claims. The existing latent-variable discussion also does not settle
this distinct identification question.

The table records the original acquisition boundaries. ``Unresolved'' describes
the evidence available to this survey. It does not mean the vendor lacks the
capability. Price and edition entitlements have not been verified for the
proposed embedding scenario \citep{exec13_ci_scope}.

\begin{longtable}{P{0.17\textwidth}P{0.34\textwidth}P{0.39\textwidth}}
\toprule
Family & Available comparison & Specific remaining source or decision \\
\midrule
\endhead
BayesFusion & Preserved structure, pattern and license documents; new causal route unresolved. & Exact intervention/query semantics and wrapper/edition support. \\
HUGIN & The initial request failed; the later GUI 9.8 manual supplies decision/information and model-class boundaries. & General identification/query semantics and the selected API/edition remain to be established. \\
Netica & Preserved learning documentation concerns adjacent operations. & Intervention and counterfactual semantics, method boundaries and license. \\
SAS & Treatment/mediation PDF transfers were incomplete; replacement HTML routes failed. & Authentic CAUSALTRT/CAUSALMED technical text and target-specific assumptions. \\
IBM & Preserved temporal-model material does not answer general identification. & Edition-specific causal/Amos methods; distinguish SEM fitting from causal identification. \\
MathWorks & Initial access was denied; later R2026b panel-method pages provide an adjacent estimation comparison. & Treatment/identification documentation with target, assumptions and runtime scope remains distinct. \\
causaLens & The documentation endpoint redirected to login. & Public technical material or subsequently authorized copies; no inference from the login page. \\
EViews & Earlier structural-model navigation supplies limited adjacent context. & Actual restrictions, intervention/scenario methods and API contract. \\
OxMetrics & PcGive/PcGets/Autometrics leads remain incompletely resolved. & Versioned identification, intervention and uncertainty methods. \\
LISREL & Retained as a structural-model candidate. & Primary model restrictions, identification diagnostics and embedding rights. \\
JMP & Earlier routes failed; retained treatment documentation is developed below. & Preserve target and edition; installed automation and rights remain unresolved. \\
Analytica & Two attempted official routes returned missing-page responses. & Influence-diagram and Decision Engine semantics, separating supplied decisions from identified effects. \\
\bottomrule
\end{longtable}

The following studies resolve parts of this initial inventory: JMP receives
a developed treatment comparison, and GeNIe/SMILE receives a supplied-model
decision and engine comparison. Those findings do not establish unrestricted
graphical identification. Chapter~\ref{scisyn:software} separately records a
named causaLens service offering; its advertised scope does not resolve the
versioned identification contract sought here. Section~\ref{gap29:software-followup}
also develops the later HUGIN decision-model and MATLAB panel-method evidence.
HUGIN's information links and discrete chance-node restriction matter to a
servicing comparison; MATLAB's panel transformations and subject covariance
matter to a branch study. Neither finding establishes a general identification
engine. Their retrieval after the literature cutoff is retained explicitly.
\citep{gap29hugin,gap29matlab}
The remaining omissions make
complete commercial coverage premature, while allowing a useful proposal
recommendation. The immediate recommendation can
use the three developed comparisons, give concrete acquisition actions for
the others, and select public exact fixtures for the next prototype. A bank
experiment is unnecessary merely to establish what a manual says. A paid trial
is also unnecessary for the mathematical reference examples below.

\section{Exact lending examples that separate the required capabilities}

The following are project constructions, evaluated with exact rational
arithmetic. They are neither fitted vendor outputs nor bank observations.
Their purpose is to make identification, partial identification and
counterfactual semantics reviewable before an expensive experiment
\citep{exec13_ci_examples}.

\subsection{The same observed loan table can support different effects}

Consider an observed law
\[
\begin{array}{c|cc}
 &Y=0&Y=1\\\hline
A=0&3/8&1/8\\
A=1&1/8&3/8
\end{array}
\]
with all four cells positive. The observed repayment difference is
$3/4-1/4=1/2$. One generating model has a binary pre-existing repayment type
$U$ with probability $1/2$, an independent bit $V$ with probability $1/4$,
offer $A=U\mathbin{\mathrm{xor}}V$, and outcome $Y=U$. Changing the offer
does not change repayment in this model: $\tau=0$.

A second model randomizes $A$ with probability $1/2$. With an independent
uniform four-state variable $W\in\{0,1,2,3\}$, let
$Y(0)=\mathbf 1\{W<1\}$ and $Y(1)=\mathbf 1\{W<3\}$.
Its effect is $\tau=1/2$, yet its observed table is identical. Enumerating
both finite models verifies the equality exactly. Additional rows from the
same observed law cannot choose between them. Observed positivity does not
resolve hidden confounding.

If the analyst commits to the second model's randomized assignment, a
treatment-effect estimate has a causal interpretation. If the information
allows both models, the same calculation is an association with unresolved
causal meaning. A commercial adapter and the local service should preserve
that difference even when they display the same number.

\subsection{Sharp bounds remain useful information}

With only consistency and binary potential outcomes, the observed successful
treated borrowers contribute $3/8$ to $\mathbb E[Y(1)]$. The unobserved
$Y(1)$ values among the untreated can contribute anything from zero to $1/2$.
Similarly, observed successful untreated borrowers contribute $1/8$ to
$\mathbb E[Y(0)]$. Thus
\[
 \mathbb E[Y(1)]\in[3/8,7/8],\qquad
 \mathbb E[Y(0)]\in[1/8,5/8],\qquad
 \tau\in[-1/4,3/4].
\]
These bounds are jointly attainable. Allocate each of the four observed
$(A,Y)$ cells to either value of its missing potential outcome. The $2^4=16$
resulting vertices preserve the table and attain both effect endpoints.
Every fractional allocation is a convex combination of the vertices, so a
linear effect cannot exceed these extrema. The exact program checks this
enumeration separately from the missing-outcome algebra.

For a hypothetical planning value of USD~1,000 per incremental repayment and
a USD~20 cost per offer, net value is $1000\tau-20$. The two equal-observed-law
models imply respectively minus USD~20 and plus USD~480 per offer. The
consistency-only range is minus USD~270 to plus USD~730. This is an assumed
value mapping, not a measured ROI or a full loan-accounting model. Its purpose
is to show why the unresolved scientific distinction can reverse a decision.

The range includes both profitable and unprofitable offers. It can justify
considering a better information design, but it does not by itself establish
that any particular experiment is worth its cost. That further argument needs
the experiment's expected resolution, the affected decisions and their costs.

\subsection{Hidden confounding does not always prevent identification}

Now assume the precise structural graph $U\to A$, $U\to Y$,
$A\to M\to Y$, with no additional arrows and independent disturbances.
Let
\begin{align*}
 \Pr(U=1)&=1/2,\\
 \Pr(A=1\mid U=u)&=1/4+u/2,\\
 \Pr(M=1\mid A=a)&=1/5+3a/5,\\
 \Pr(Y=1\mid M=m,U=u)&=1/10+2m/5+2u/5.
\end{align*}
The mediator intercepts all directed effects of $A$ on $Y$ in this stipulated
law. The observed functional
\[
 \mu_a=\sum_m P(m\mid a)\sum_b P(Y=1\mid m,b)P(b)
\]
equals the intervention mean here. To verify rather than assume this equality,
note that $U$ and $M$ are conditionally independent given $A$, so
\[
 P(Y=1\mid m,b)=1/5+2m/5+b/5.
\]
Averaging over $P(A=b)=1/2$ gives $3/10+2m/5$; averaging over
$P(M=1\mid A=a)$ gives $19/50+6a/25$.
Independently summing the intervened structural law over $(U,M)$ produces
the same two means, $19/50$ and $31/50$. The effect is $6/25=0.24$,
while the ordinary observed difference is $11/25=0.44$.

The example is a useful front-door fixture, not a guarantee for every mediator.
It would require another argument if the offer had a direct effect on repayment,
if an unmeasured cause affected both offer and mediator, or if the measured
mediator failed to represent the stipulated variable. The immediate software
requirement is nonetheless clear: a failed observed back-door adjustment must
not be reported as a universal nonidentification proof. Under the hypothetical
USD~1,000/20 mapping, this model's identified net value is USD~220 per offer;
the model assumptions carry that interpretation.

\subsection{Even complete randomized margins can leave an individual question open}

Let $A$ and $U$ be independent fair bits. In one world set
$Y(0)=Y(1)=U$; in a second set $Y(0)=U$ and $Y(1)=1-U$.
Both worlds have $P\{Y(a)=1\}=1/2$ for each action, the same randomized
observed law, and average effect zero. In the first world, a treated borrower
who repaid would also have repaid without the offer, so $q=0$. In the second,
that borrower would not have repaid, so $q=1$.

The same margins allow the population probability of benefit
$P\{Y(1)=1,Y(0)=0\}$ to range from zero to $1/2$. A randomization result
can identify an average effect while leaving this cross-world association
unresolved. An abduction engine supplied with one complete structural model
may return an answer, but the evidence must say which coupling it used. A
language model's plausible borrower narrative is not additional identifying
information.

\section{A proposal architecture that can be compared fairly}

The smallest coherent prototype should expose a query object containing the
population, outcome clock, intervention, conditioning event and target units.
It should attach the supplied graph or structural model, observed-variable
map, permitted data laws and assumptions. The method registry should state
the exact domain of each method. A response should then carry its expression
or set, derivation/certificate where available, unresolved premises and the
next permissible operation.

Estimation belongs after an identified functional has been selected. A
licensed Stata adapter, a documented graphical intervention adapter, or an
open-source estimator could occupy that stage. The comparison should hold the
query, data, adjustment set and sampling assumptions fixed. Method-specific
diagnostics, missing values and failed requests remain part of the result.
Numerical agreement is meaningful only when both tools answered the same
question under the same assumptions.

For individual counterfactuals, a separate structural-model path should retain
the posterior exogenous distribution after abduction and the exact action used
in prediction. It should refuse to imply unique individual identification when
the available observations permit several couplings. For partial identification,
the interface should preserve closed/open endpoints, unions and empty sets
where the method requires them. These are proposed requirements; the local
back-door reference implementation and deterministic workflow checks establish
only their documented subset.

The exact fixtures above permit a useful next experiment with no authentic
bank data. First compare each proposed method's supported-domain response to
the independently derived truth: equal observed laws, a front-door functional,
sharp bounds and differing cross-world couplings. Then add finite public or
generated samples under a frozen observation law, with uncertainty, availability
and width criteria declared before final draws. A vendor runtime comparison
can follow once technical contracts and licensing are resolved. Neither an
unavailable commercial manual nor a deferred paid trial needs to block the
public prototype route.

The investment case should consequently distinguish three decisions: adopt
an existing estimator/workflow where its domain and rights match; build the
missing scientific contract and cross-tool explanation where that provides
demonstrable value; and defer claims requiring undocumented or unidentified
capabilities. This makes the commercial landscape consequential for the
proposal while retaining honest limits on source coverage and implementation.
The next study makes the identification and abduction contracts more concrete;
the subsequent estimator and product comparisons retain the same distinction
between a defined target, an implemented operation and justified uncertainty.
\label{exec13:commercial-identification-end}


% SOURCE docs/release/2026-10-29/source/docs/writing-repair/2026-10-29/src/bb6773e501a4fb48.tex CHAR 0
\section{Edition and integration choices for the lending query}
\label{exec14:commercial-contracts}

An actionable comparison must connect the banking request to the particular
operation that an edition exposes. The population effect of a repayment offer,
the effect transmitted through payment behavior, and an alternative outcome
for an already observed borrower should therefore have different request
contracts. More detailed Bayes Server documentation supplies a useful test of
this distinction. Its Advanced Causal Query guide describes back-door and
front-door criteria, subsequent adjustment calculations, multiple treatments
and joint outcomes. Support for hidden confounders and direct effects varies
with the selected algorithm. Those are documented capabilities with which the
proposal must compete; the list is not evidence of a complete general
identification algorithm \citep{exec14_ci_causal_query}.

The version evidence has several layers. The guides describe the features as
available from version~10. The acquired release history associates advanced
causal queries, back-door and front-door algorithms, causal effects analysis
and the abduction helper with version~10.2 beta, dated 19 June 2022. Its highest
listed release is 11.4.0.0, dated 17 April 2026. These observations support a
version-family comparison; they do not identify the bytes of an installed
assembly or establish equivalence across its interfaces
\citep{exec14_ci_release_history}. The copies were acquired on 14 October
2026. As elsewhere in this survey, dated release entries and copyright notices
do not authenticate a mutable page's exact contents on the 7 October cutoff.

\subsection{A narrow identification adapter has a concrete contract}

The documented \texttt{FrontDoorCriterion} class takes a supplied network. Its
\texttt{Identify} methods receive interventions, outcomes, other evidence
and options, and return criterion-specific identification information without
performing estimation. Associated methods, \texttt{IdentifyXZ} and
\texttt{IdentifyZY}, address the adjustment subproblems between treatment,
front-door nodes and outcomes. The documentation explicitly envisages using
the identified node sets outside Bayes Server. It also distinguishes a Boolean
validity check from validation that raises an exception
\citep{exec14_ci_api_frontdoor}. This is a much more useful integration lead
than a claim that a product ``supports causality.''

For the repayment example, a proposed adapter would receive the offer variable
$A$, a mediator $M$, repayment $Y$, an observed-variable map, a supplied causal
network and the desired intervention. It would return the criterion, the
identified node sets, the input model identity and the query to which they
apply. An external estimator would then require the observed distributions
and support conditions needed to evaluate the functional. The exact
front-door calculation developed earlier illustrates that separation:
returning a mediator name is not the same as fitting the required conditional
distributions, and fitting them is not proof that the graph describes the bank.

A useful minimal adapter could therefore export a reviewed identification
result before undertaking a much larger attempt to embed an entire modeling
environment. The export must preserve node identity, intervention versus
ordinary evidence, treatment coding, relevant clock and units. Its response
should distinguish a valid criterion result, invalid input, unsupported
options and failure to find a set. In particular, failure of this one
criterion must not become a universal nonidentification certificate. That
conclusion would need a complete algorithm for the declared graph/query class
or an independent argument. The selected documentation establishes neither
general ID/IDC completeness nor unrestricted counterfactual identification.

The reference revision examined in this comparison exposed supplied-DAG
back-door checks, supplied-moment Gaussian residual inversion and explicit
sensitivity and clock operations. It did not expose this front-door adapter,
a general abduction engine or a fitted mediation service. Later fitted-summary
inference in Chapter~\ref{gap29:public} is separate evidence and does not
qualify those missing commercial interfaces. The proposed
scientific contract can organize those future capabilities, but it does not
erase the present functional gap. An investment case should demonstrate the
value of preserving scientific meaning across an existing estimator and a
specific identification service before proposing another broad implementation. \citep{exec14_ci_comparison}

\subsection{Effect labels need mathematical definitions}

The Advanced Causal Query guide describes total and direct effects using paths.
One illustrative total-effect list includes a path of the form
$A\leftarrow C\rightarrow Y$. That is a back-door association path, not a
directed causal route from $A$ to $Y$. This wording should not be adopted as
the definition of total effect, and it does not by itself demonstrate an
implementation error \citep{exec14_ci_causal_query}. A wrapper must bind the
requested quantity to an explicit mathematical expression and then determine
whether the selected vendor operation computes it.

For example, the total contrast is
$\mathbb E[Y(1)-Y(0)]$. A controlled direct contrast at mediator level $m$ is
$\mathbb E[Y(1,m)-Y(0,m)]$, while a natural direct contrast uses the
unit-specific mediator $M(0)$ inside a nested potential outcome. A menu label
about direct links does not settle which of these interventions and
cross-world premises is intended. Stata's specified mediation workflow
provides a useful contrast because it states its natural-effect targets and
confounding restrictions \citep{exec13_ci_stata_mediate}. The proposal should
compare matching estimands, not merely similarly named report columns.

\section{Abduction must preserve uncertainty about the borrower}

The Bayes Server abduction dialog describes querying characteristic variables
given factual evidence, removing evidence for the individual and setting
evidence on those characteristics. Its .NET \texttt{Abduction.Update} signature
separates current evidence, the variables supplying abduction evidence, the
characteristic variables and additional options. This is explicit support for
an abduction step. The selected pages do not fully specify how an uncertain
multivariate posterior is represented or propagated through every supported
case \citep{exec14_ci_abduction,exec14_ci_api_abduction}. That remaining
question matters even when the conditional mean of each characteristic can
be computed accurately.

Let $U$ represent latent repayment capacity in normalized units. For an
illustrative required payment burden $a$, define
\[
 U\sim\operatorname{Uniform}\{0,1,3\},\qquad
 Y(a)=\mathbf 1\{U\geq a\}.
\]
Assignment of the burden is independent of $U$ in this stipulated model.
For a borrower observed at $A=1$ with repayment $Y=1$, Bayes' rule eliminates
$U=0$ and retains $U=1$ and $U=3$ with equal posterior mass. Under the
alternative burden $a'=2$, the individual-conditional query is
\[
 \Pr\{Y(2)=1\mid A=1,Y=1\}
 =\sum_{u\in\{1,3\}}\mathbf 1\{u\geq2\}
   \Pr(U=u\mid A=1,Y=1)=\frac12.
\]
The population intervention probability is instead $1/3$. Replacing the
posterior of $U$ by its mean, $\mathbb E[U\mid A=1,Y=1]=2$, gives the
indicator $\mathbf 1\{2\geq2\}=1$. All three calculations concern the
same structural equation, but they answer different questions or use a lossy
approximation. The finite enumeration reproduces the fractions exactly; it
is a project construction, not a test of Bayes Server
\citep{exec14_ci_exact}.

More generally, a model-specified individual counterfactual averages the
alternative structural response over the posterior of the same exogenous
state that generated the factual observations. Applying a nonlinear response
to a posterior mean need not equal that average. For several exogenous
variables, preserving only separate marginal distributions can also discard
dependence needed by the response. A capability contract should state whether
the operation returns a uniquely recovered characteristic value, a posterior
distribution, a point approximation or a set of compatible states. It must
also identify the extra structural assumptions that connect alternative
outcomes for the same borrower.

To make the financial consequence explicit, stipulate a USD~1,000 payment
upon success and a USD~20 cost under the alternative scenario. The posterior
calculation gives an expected payoff of $1{,}000(1/2)-20=480$ dollars. The
mean substitution gives $1{,}000(1)-20=980$ dollars, an overstatement of
USD~500 in this model. These are alternative-scenario payoffs, not incremental
policy profit: a decision comparison would additionally specify the baseline
payoff and all costs. Neither the amounts nor the latent law are bank
measurements \citep{exec14_ci_exact}.

The immediate public/generated test is therefore precise. Supply a small
structural model whose exogenous posterior has more than one possible state,
compare its exact enumerated counterfactual distribution with the returned
object, and repeat with dependent exogenous variables. Test removal of factual
evidence separately from the intervention. Retain unsupported or approximate
behavior explicitly. The inspected signature is sufficient to motivate this
prototype experiment; it is insufficient to declare the result in advance.

\section{Sweeps, units and uncertainty in the commercial output}

The Effects Analysis guide describes sweeps over intervention values. Discrete
treatments use their states; continuous treatments are discretized conditional
on current evidence. For discrete outcomes the requested state probability is
reported; continuous outcomes can produce a mean and variance. The guide also
normalizes the chart's horizontal axes when treatments are compared
\citep{exec14_ci_effects}. Those choices affect what an exported plot means.

A lending manager should be able to recover the actual payment burden or
interest-rate change corresponding to every displayed point. A normalized
coordinate without the original dose, unit and transformation is not an
actionable pricing instruction. If the grid depends on factual evidence, a
different borrower's curve may use different underlying values. A prospective
adapter must carry the dose grid and conditioning evidence as data, rather
than reconstructing them from an image. The intervention guide further
describes discrete, continuous and dynamic-network interventions, which must
remain distinct from ordinary observational evidence
\citep{exec14_ci_interventions}.

The reported variance of an outcome under a fixed fitted model is also a
different object from uncertainty about its estimated causal mean. If
$\theta$ denotes parameters fitted from data, then
$\operatorname{Var}_{\theta}\{Y(a)\}$ describes outcome variability at
that model, while sampling variation of an estimator of
$\mathbb E_{\theta}\{Y(a)\}$ requires an estimation and sampling law.
Uncertainty about the graph or about an unmeasured causal premise is another
layer. A mean/variance pair from an intervention sweep cannot be relabeled as
a calibrated confidence interval for the offer effect. The proposed result
format should identify these layers and report which are included.

\section{Reuse choices across the compared families}

BayesiaLab's supplied-graph and data-mapping workflow remains a credible
commercial comparator. Its Engine API overview distinguishes Java modeling,
learning and inference configurations from REST inference. The table places
Java model files on customer servers and REST model files on Bayesia servers
in France; production capacity is tied to permitted concurrent instances,
with documented license-validation connectivity. These facts matter to an
integration design, but they do not establish that every graphical causal
function is exposed by every configuration
\citep{exec13_bayesia_api,exec13_bayesia_professional}. The proposed route is
a narrowly specified graph/query export followed by verified interpretation
of the returned result. The missing causal API mapping must be acquired before
committing to an inference-only endpoint for a more elaborate task.

Stata supplies a mature treatment and mediation baseline for specified targets.
Its treatment-effect documentation is more directly relevant to a one-time
offer ATE or ATT than a generic graph-learning comparison
\citep{exec13_ci_stata_causal_intro}. The license-options text distinguishes
an identified user's installations from concurrent network use at a physical
site. This is useful deployment context, but the options page does not settle
the entitlement for a proposed multiuser service
\citep{exec13_ivrd_license}. A first comparison can use an appropriately
authorized analyst workflow and complete output artifacts. Any subsequent
CLI/MCP service needs a separately resolved execution and export contract.

The retained JMP 19.2 documentation supplies the treatment comparison that
the initial acquisition had left unresolved. The JMP Pro Causal Treatment personality describes ATE and
ATET for binary treatment, and potential-outcome means for continuous treatment.
Its assumptions distinguish consistency, positivity and conditional
exchangeability. Its binary estimators use specified propensity and outcome
models, with weighting, regression adjustment and augmentation alternatives
\citep{exec13_ch_jmp_causal}. These are meaningful commercial capabilities
that the local reference runtime does not yet provide.

The documentation's potential-outcome predictions should not be described as
identified missing individual outcomes. Model-based conditional means and
population averages have their own legitimate roles; they do not supply the
cross-world dependence in the borrower example. A proposed file-based
comparison should preserve target weights, observation exclusions, nuisance
specifications and treatment coding before attempting unattended automation.
Only inspected prose supports this comparison; image-only equations were not
treated as inspected technical evidence. Installed build, automation rights
and selected deployment prices remain unresolved.

BayesFusion's SMILE documentation provides a compiled C++ library and named
wrappers, as well as separate development and deployment licensing categories
\citep{exec13_smile_intro,exec13_smile_license}. That is adjacent integration
evidence. Section~\ref{sec:exec24-smile-reuse} develops the supplied-model
decision and temporal interfaces; a matched general intervention/counterfactual
contract remains a separate qualification.
For HUGIN, Netica, SAS, IBM, Mplus, MathWorks, causaLens, EViews, OxMetrics,
LISREL and Analytica, the technical or edition-specific questions retained
earlier in this chapter remain open at different levels. Some have useful
adjacent manuals or original-paper leads; inaccessible or incomplete pages
do not establish absent functionality. Broader ranking should wait for the
matching source rather than borrow a conclusion from another software area. \citep{exec14_ci_comparison}

\section{Published prices and conditions for a justified comparison}

The acquired Bayes Server purchase page explicitly selects Commercial and GBP
in its form. It lists version-11 Professional individual UI and API developer
licenses at GBP~550 each, and their combination at GBP~1,100, before applicable
VAT. It describes perpetual licenses for the major version with minor upgrades
and per-developer API functionality within the purchaser's own product
\citep{exec14_ci_purchase}. These are public page observations on
14 October 2026, not a quotation, a historical-cutoff price or total bank
deployment cost. They should remain in their original currency; conversion
would need an explicit exchange-rate date and cost purpose.

The linked, undated software agreement contains restricted executable-use and
transfer terms and a requirement for prior written consent before disclosing
benchmark results to third parties. It allows that an amendment or addendum
may accompany the software \citep{exec14_ci_eula}. This does not justify a
blanket conclusion that the purchase-page description is false or that a
particular proposed distribution is licensed. It identifies an exact remaining
commercial question: reconcile the selected order and applicable agreement
before promising embedded deployment or public benchmark publication. No
license was accepted, no vendor was contacted and no commercial binary was
executed for this document comparison.

The proposal can proceed with a useful staged decision. First define a narrow
target and compare its documented contract, information requirements and
failure states against the local service. Next build a generated/public-data
prototype that demonstrates the missing semantic preservation and justifies
the cost of a matched experiment. Where a proprietary engine is involved,
resolve the actual execution and reporting terms before claiming a measured
comparison. Finally measure accepted output quality, review time, engineering
effort and operating costs together. A lower license amount, additional menu
entry or successful API response cannot alone decide which route gives the
bank a more defensible answer.
\label{exec14:commercial-contracts-end}


% SOURCE docs/release/2026-10-29/source/docs/writing-repair/2026-10-29/src/6768426fe839ace4.tex CHAR 0
\section{Commercial heterogeneity tools: from a conditional effect to a servicing decision}
\label{exec13:commercial-heterogeneity}

A servicing manager deciding whom to contact needs more than a ranking of
predicted recoveries. An account may have a high probability of repayment
without a call, a large incremental response to a call, or a large response
whose delivery cost exceeds its value. These are different quantities.
Commercial comparison must therefore start with the population, intervention,
outcome horizon and decision rule, then ask which documented estimator matches
that question. A label such as individualized effect cannot establish that
both potential outcomes of a particular customer have become observable.
The Stata manual itself defines its individualized average treatment effect
as an average for individuals sharing the specified characteristics
\cite{exec13_ch_stata}.

The comparison here concerns a declared universe of seventeen commercial
families. Complete primary documents support a detailed comparison of Stata
and two distinct JMP pathways, together with a more limited DataRobot
weighting workflow. Saved technical material was acquired on 13 October 2026.
Its version and publication cues precede the 7 October literature cutoff,
but those cues do not authenticate every byte as it appeared on the cutoff
date. The incomplete family coverage, failed requests and unreviewed rights
are substantive limits, not evidence that a competitor lacks a feature.
No commercial product was executed for this comparison
\cite{exec13_ch_scope}.

\subsection{One population, several legitimate answers}

Let $D$ denote a binary servicing action, $Y(d)$ the recovery in dollars over
one fixed horizon under action $d$, and $G$ a risk group measured before the
action. In the generated example, lower-risk accounts have population mass
$3/4$, action probability $1/10$, effect $2$ dollars and delivery cost $1$
dollar. Higher-risk accounts have mass $1/4$, action probability $4/5$, effect
$10$ dollars and cost $12$ dollars. Potential outcomes are stipulated, with
$Y(0)=8$ and $4$ in the respective groups and $Y(1)=Y(0)+\tau_G$. Treatment
is randomized within each group. Thus conditional exchangeability and
positivity hold by construction; they are not conclusions drawn from a bank
dataset. The exact finite-population review uses thirty lower-risk and ten
higher-risk accounts, with three and eight treated respectively. This
representation realizes the declared probabilities exactly
\cite{exec13_ch_exact}.

The average effect in the eligible population is
\[
 \tau_{\mathrm{ATE}}=\tfrac34(2)+\tfrac14(10)=4.
\]
The treated fraction is $11/40$, and the average effect among treated
accounts is instead
\[
 \tau_{\mathrm{ATET}}
 =\frac{(3/4)(1/10)(2)+(1/4)(4/5)(10)}{11/40}
 =\frac{86}{11}.
\]
The larger second number reflects the concentration of treatment in the
higher-effect group. Reporting it as the gain from treating every eligible
account would change the target population. Conversely, the ATE does not
answer a question restricted to the accounts currently receiving treatment.
These differences require a target label in an interchange record, even if
two interfaces both expose a field called effect
\cite{exec13_ch_exact}.

Residualization introduces a further distinction when effect modification
is omitted from the conditioning set. Write $W$ for all nuisance controls,
$e(W)=\Pr(D=1\mid W)$ and $h(W)=\E[Y\mid W]$. If the requested effect
model is constant while the actual effect varies with $W$, the population
least-squares residual slope is
\begin{align*}
 b_{\mathrm{PO}}
 &=\frac{\E[(D-e(W))(Y-h(W))]}{\E[(D-e(W))^2]}\\
 &=\frac{\E[e(W)(1-e(W))\tau(W)]}{\E[e(W)(1-e(W))]}.
\end{align*}
The second equality follows by conditioning on $W$, using the Bernoulli
variance of $D$ and the stipulated conditional treatment law. It is an
overlap-weighted effect, which is $214/43$ in this example, rather than the
ATE of $4$. This is a target-under-misspecification calculation, not a
claim that the inspected Stata implementation silently promises an
unrestricted ATE. Its partialing-out derivation specifies an effect
$\tau(x)$ that does not additionally vary with nuisance-only controls $w$.
Changing that scientific restriction changes the interpretation of the
residual moment \cite{exec13_ch_stata,exec13_ch_exact}.

For comparison, define $\mu_d(W)=\E[Y(d)\mid W]$ and the oracle augmented
inverse-probability score
\begin{align*}
 \psi={}&\mu_1(W)-\mu_0(W)\\
 &+\frac{D\{Y-\mu_1(W)\}}{e(W)}
 -\frac{(1-D)\{Y-\mu_0(W)\}}{1-e(W)}.
\end{align*}
With the true conditional functions and exchangeability, each residual
term has conditional expectation zero. Hence $\E[\psi]=\E[\tau(W)]=4$.
Conditioning the same score on a coarser pre-action characteristic $X$
gives $\E[\psi\mid X]=\E[\tau(W)\mid X]$ when $X$ is included in the
conditioning information. Estimating those functions, fitting a conditional
mean and obtaining valid uncertainty require additional regularity and
sampling arguments. The exact score calculation alone establishes neither
finite-sample calibration nor robustness to arbitrary learner errors
\cite{exec13_ch_exact}.

This example motivates a concrete proposal requirement: record the eligible
population, treated subpopulation, effect modifiers, nuisance-only variables
and weighting target separately. A coarse risk-band request must not inherit
an estimand merely because the nuisance model uses many variables. Model
specification, target aggregation and source-population weights are part of
the scientific contract. They must survive export to a library result, a CLI
artifact and an MCP response.

\subsection{Stata: a developed conditional-effect workflow}

The inspected \texttt{cate} documentation describes continuous outcomes and
binary treatment, with partialing-out and augmented inverse-probability
weighting paths. Outcome and treatment nuisance functions can use lasso,
random forests or documented parametric alternatives; the conditional effect
can use a forest or linear regression. It distinguishes IATEs, prespecified
group average effects, and sorted groups constructed from estimated IATE
quantiles. The distinction matters operationally: a regulator-defined risk
band is a different grouping mechanism from a band discovered after fitting
the same sample \cite{exec13_ch_stata}.

The default cross-fitting specification randomly divides observations into
ten folds. Default honest forests split the tree-building sample between
split selection and leaf estimation. The manual explains that disabling
honesty disables the corresponding forest standard errors and confidence
intervals. These are useful documented controls, but their purposes should
remain distinct. Cross-fitting separates nuisance estimation from the
observation scored; honesty separates tree partitioning from leaf labeling.
Neither operation, by its name alone, demonstrates validity when accounts
repeat within households or outcomes share calendar shocks. The inspected
random-observation algorithm does not establish a household- or
time-dependent inference result; that comparison remains open rather than
being converted into a claim about every Stata command
\cite{exec13_ch_stata}.

The out-of-bag alternative requires forest outcome and treatment models.
It cannot be combined with the documented sorted-group option, and it is
an alternative to the cross-fitting setting. A proposed wrapper must
validate these combinations before launching a job. It should not expose
cross-fitting, out-of-bag prediction and arbitrary grouping as freely
composable switches. The default propensity tolerance is $10^{-5}$; the
manual describes an error for estimated propensities below that tolerance
and an option identifying offending observations. This is a documented
computational safeguard, not a certificate of scientific positivity
\cite{exec13_ch_stata}.

Missingness is similarly method-specific. The documentation permits missing
covariates when all three models use forests, whereas other configurations
omit affected observations; an explicit option forces listwise omission.
Changing learner settings can therefore change the analyzed population.
For servicing, the disappearance of accounts with incomplete contact history
may alter both estimated response and the operational population. A wrapper
needs input and retained-row counts, reasons for exclusions, and a clear
distinction between a model's missing-value routing and a causal assumption
about the missingness process \cite{exec13_ch_stata}.

Postestimation includes conditional predictions, heterogeneity tests, group
comparisons, projections and evaluation of a prespecified treatment rule.
Forest prediction standard errors use a bootstrap of little bags; the
linear-regression prediction path uses a delta method. The policy command
computes the value of a specified rule and can compare two rules. These
are concrete comparison requirements for the proposed library: a point
estimate without a target, prediction method, uncertainty convention and
policy-selection history would be a weaker deliverable. At the same time,
documentation of a prespecified-policy command is not proof that a policy
chosen to maximize the same evaluation outcomes has valid uncertainty
\cite{exec13_ch_stata}.

Stata is therefore a candidate for a licensed external reference execution,
subject to edition, automation and runtime rights. The proposal should seek
matched generated-data comparisons for ATE, conditional means and
prespecified policy value, with identical samples and effect modifiers.
It should not propose copying proprietary implementation code or equate
a CLI wrapper with validated cross-language statistical equivalence.
Exact version, seeds, defaults, retained sample and inference method would
need to accompany each result. Runtime and entitlement verification remain
open in this survey \cite{exec13_ch_scope}.

\subsection{JMP Uplift and JMP Pro Causal Treatment answer different questions}

The JMP Uplift overview describes partition models whose splits maximize
treatment differences. For a continuous outcome, candidate splits enter a
linear model through split, treatment and their interaction. For a
categorical response, the corresponding comparison uses a logistic model.
Predictor selection uses interaction significance with a Monte Carlo
adjustment for the number of candidate levels or models. This documented
adjustment addresses a predictor-selection mechanism; it is not a general
proof of causal identification or uncertainty after all possible model and
policy searches \cite{exec13_ch_jmp_uplift}.

The separately versioned JMP 16 guide makes two particularly consequential
interface details explicit. Treatment value ordering determines the
comparison, and more than two treatment levels are reduced to the first
level against all remaining levels combined. Validation data help choose
the split sequence. Those older-guide details must be rechecked in a
chosen contemporary runtime; they are not automatically a JMP 19 guarantee.
Their architectural lesson already applies: an adapter accepting three
servicing actions must not silently turn a three-action decision into a
binary first-versus-rest contrast. Likewise, validation used for selecting
a tree is not a fresh final policy-evaluation sample. A node's displayed
treatment difference or test statistic does not by itself settle that
selection issue \cite{exec13_ch_jmp_uplift}.

JMP Pro's Causal Treatment personality is a different pathway. The inspected
19.2 documentation describes binary or continuous treatment, binary or
continuous outcome, and average-effect summaries. For binary treatment,
the report distinguishes population ATE from ATET and documents
inverse-probability weighting with ratio adjustment, regression adjustment
and augmented weighting. The specified propensity model is logistic
regression, with linear or logistic outcome regression according to
outcome scale. These claims concern those documented model paths, not an
inferred native conditional-effect forest or arbitrary machine-learning
interface \cite{exec13_ch_jmp_causal}.

For continuous treatment, the prose specifies a normal conditional
treatment density with parameters fitted by linear regression. It presents
a stabilized-weight marginal structural model and a generalized propensity
score route to mean potential outcomes at specified doses. Thus a continuous
contact intensity or price change needs a dose and support definition,
not a reinterpretation of a binary-treatment coefficient. The displayed
report discusses a treatment and squared-treatment regression; its
image-rendered equations were not transcribed into this comparison.
Only inspected prose supports the implementation description
\cite{exec13_ch_jmp_causal}.

The same documentation states consistency, conditional exchangeability and
positivity as identification assumptions. Empirical overlap plots can
reveal severe support problems in observed data; they cannot establish the
absence of unmeasured confounding, a well-defined intervention or population
positivity everywhere. An attractive diagnostic plot must therefore be
stored with its interpretation and limitations. The proposal should
compare JMP's documented reports as a user-facing baseline while retaining
explicit assumption and target records \cite{exec13_ch_jmp_causal}.

The practical reuse choices differ by pathway. Uplift is a candidate
benchmark for transparent treatment-response segmentation under a suitable
experimental design. Causal Treatment is a candidate comparator for the
specified average and dose-response models. Neither choice establishes
availability of an unattended server interface, MCP service or permitted
redistribution. Product edition and rights must be reviewed before an
execution adapter is approved. A file-based result import can be designed
now, with an explicit unresolved-runtime disposition
\cite{exec13_ch_scope}.

\subsection{A documented platform workflow is a narrower claim than a native estimator}

DataRobot's official Causal AI accelerator overview describes fitting a
treatment-propensity model, calculating inverse-probability weights and
comparing outcomes in the weighted population. Its example concerns
hospital readmission; the relevant architectural point is the separation
between predictive propensity fitting and causal weighting. The overview links to a notebook whose pinned source is developed in
Section~\ref{sec:exec25-datarobot-reuse}. That study traces in-sample propensity
predictions, row/class mapping and normalized weighting, while retaining the
unresolved source-permission and native-runtime boundaries. Neither the
overview nor that source inspection establishes a native CATE engine, a
doubly robust score or calibrated inference
\cite{exec13_ch_datarobot}.

For the proposed library, the conditional adoption route is to accept
versioned nuisance predictions through a narrow data contract if their
training population, folds, target encoding and preprocessing are known.
Identification, overlap policy, target aggregation and inference remain
separate responsibilities. A prediction service becomes usable in an
orthogonal-score workflow only when its training/evaluation separation
and error conditions are justified. The current evidence does not approve
that integration or a performance comparison. It specifies what code and
rights review must precede one.

The other fourteen families are not silently dropped. BayesFusion, Bayesia,
HUGIN, Norsys, IBM, Mplus, EViews, OxMetrics, Scientific Software
International, Bayes Server and Lumina retain prior-area source leads,
but a matched heterogeneous-effect method, edition and rights record is
not established by those leads. SAS's newly saved treatment-effects PDF
is partial and excluded; failed technical routes do not settle the scope
of its products. The requested MathWorks route returned access denial,
and the causaLens request failed certificate-hostname validation. These
are acquisition limits. The survey requires successor source work before
ranking those families against the proposed capability
\cite{exec13_ch_scope}.

\subsection{Overlap, selection and the economics of a policy}

Overlap policy is not a minor numerical option. Consider a generated
population with true propensity $1/20$, $Y(1)=10$ and $Y(0)=0$. If an
unnormalized IPW expression replaces the propensity by a floor of $1/10$,
its expected treated contribution is $(1/20)10/(1/10)=5$, although the
true ATE is $10$. This deliberately narrow counterexample concerns that
unnormalized formula. It does not assert the same bias for a ratio-normalized
estimator, a correctly specified augmented estimator or a separately
defined overlap target. It shows why the adapter must preserve the exact
weighting and normalization rule, not just a boolean clipping flag
\cite{exec13_ch_exact}.

Policy value introduces costs and selection. Let $\pi(G)$ be a fixed
binary contact rule and $c(G)$ its per-contact cost. Its incremental net
value per eligible account in the constructed example is
\[
 V(\pi)=\E[\pi(G)\{\tau_G-c(G)\}].
\]
Treating no one yields zero; treating only lower-risk accounts yields
$3/4$ dollars; treating only higher-risk accounts yields $-1/2$ dollars;
treating both yields $1/4$ dollar. The highest-effect group is not the
highest-net-benefit group because its stipulated delivery cost is larger.
These are exact scenario values under fixed effects, costs and population
shares. Capacity, customer harm, uncertainty, interference and treatment
feasibility would require additional constraints before a bank decision
\cite{exec13_ch_exact}.

Even without those complications, choosing a policy on its evaluation
noise can inflate the apparent benefit. Suppose two equal true effects
are observed with independent errors taking $-1$ and $1$ with equal
probability. The expected larger reported effect is $1/2$. Selecting
with these observations and evaluating with fresh independent zero-mean
errors gives expected evaluation error zero. This elementary construction
does not accuse a vendor of using an invalid procedure. It explains why
the proposed workflow must distinguish policy training, tuning and final
evaluation, and why a prespecified-rule command cannot automatically
certify an adaptively chosen portfolio \cite{exec13_ch_exact}.

\subsection{What the proposed capability must add, reuse and leave unresolved}

The inspected alternatives make the required comparison concrete. A
credible proposed heterogeneous-effect module needs documented target
definitions and aggregation weights; nuisance and effect-model options;
reproducible splitting and missingness rules; overlap diagnostics and
failure behavior; uncertainty with its sampling-unit assumptions; and
policy-value reporting under explicit cost and selection conventions.
The commercial documentation establishes comparison requirements. It does
not demonstrate that this project's current service already implements
all of them \cite{exec13_ch_stata,exec13_ch_jmp_causal,exec13_ch_scope}.

A first implementation should use a small declared set of targets and
sampling laws. An artifact would bind the treatment contrast, outcome
horizon and currency unit; the eligible and retained populations;
pre-action effect modifiers and controls; nuisance predictions and their
fold membership; weight transformation and normalization; conditional
versus population aggregation; and the uncertainty and policy-evaluation
method. Unsupported multi-action, repeated-account or selected-policy
requests should receive a specific disposition rather than a
misinterpreted numeric result. CLI and MCP can carry this record, but
transport does not supply the absent statistical argument.

The recommended comparison sequence is to reproduce the exact target
examples, then use a prospectively specified generated design for matched
estimator and fixed-policy comparisons. Vendor execution is conditional
on a licensed, versioned runtime. Code porting requires an actual right
to reuse code; the present recommendation for proprietary methods is
documented comparison or authorized external execution. Open-source
implementation choices remain subject to their separate source, license
and numerical reviews. The immediately actionable contribution is a
scientifically explicit interface and evaluation design, with the
commercial survey still inconclusive where essential method, edition,
rights or runtime evidence is missing \cite{exec13_ch_scope}.
\label{exec13:commercial-heterogeneity-end}


% SOURCE docs/release/2026-10-29/source/docs/writing-repair/2026-10-29/src/f91385df2a3f9ea8.tex CHAR 0
\section{Commercial instruments and credit eligibility thresholds}
\label{exec13:commercial-iv-rd}

A credit committee changes an eligibility rule at a score threshold. Approval
increases above the threshold, but some eligible borrowers decline the loan,
some ineligible borrowers receive an exception, and a second benefit may
change at the same score. The committee wants the value of extending credit,
not merely the difference in outcomes between score groups. This example
connects instrumental variables (IV) and regression discontinuity (RD), while
showing why selecting a familiar command does not settle identification.
The software must represent the assignment rule, actual action, running
variable, cutoff, observation horizon and uncertainty in the local first
stage. The analyst must defend what the assignment changes besides the action.

Commercial software already provides substantial parts of this workflow.
Stata documents weak-instrument test inversion with disconnected, unbounded
and empty sets. SAS and RATS provide systems and instrument specifications;
EViews provides several distinct instrument diagnostics; SPSS exposes a
two-stage least-squares procedure; and JMP documents equation-specific
model-implied instruments. A proposal that compares only conventional least
squares with its own set-valued interface would miss these alternatives.
The relevant comparison concerns the target, assumptions, numerical boundary,
licensed execution and preservation of the result between applications.

The comparison below uses inspected official technical text for six commercial
families and pinned upstream Stata RD packages. MATLAB, PcGive and LIMDEP
remain candidates with unavailable essential originals in this bounded
review. They are not classified as lacking IV or RD. Documentation identifies
the available operation in its stated edition; it does not establish runtime
performance or a successful integration here. Mutable manuals acquired on
13 October 2026 do not prove their exact bytes existed at the 7 October
literature cutoff. Dated editions and pinned package commits narrow that
uncertainty, and unresolved edition and execution rights remain explicit
\citep{exec13_ivrd_survey}.

\subsection{A threshold jump is a target with assumptions}

Write $S$ for the score, $c$ for the eligibility threshold,
$Z=\mathbf 1\{S\ge c\}$ for eligibility, $A$ for actual credit extension
and $Y$ for subsequent value. Where the relevant limits exist and are
identified by the score support, define
\[
 \Delta_Y=\lim_{s\downarrow c}E[Y\mid S=s]
          -\lim_{s\uparrow c}E[Y\mid S=s],\qquad
 \Delta_A=\lim_{s\downarrow c}E[A\mid S=s]
          -\lim_{s\uparrow c}E[A\mid S=s].
\]
The ratio $\Delta_Y/\Delta_A$ is a local observed contrast when
$\Delta_A\ne0$. Its causal interpretation needs a model. Under continuity
of the relevant potential-outcome and compliance laws, an exclusion
restriction, a nonzero first-stage jump and local monotonicity, a binary
action can give a local complier target. Monotonicity rules out borrowers
whose treatment response moves in the opposite direction to eligibility.
It does not follow from observing a positive aggregate first stage. A
constant structural-effect model instead supplies its own stronger
interpretation. Neither target is automatically the bank-wide effect of
extending credit to every applicant.

The local complier calculation can be seen directly in a stipulated
limiting population whose potential outcomes and treatment responses
have the same distribution on both sides of the cutoff. Under exclusion
and consistency, write $A(1),A(0)$ for eligibility responses and
$Y(1),Y(0)$ for action outcomes. Monotonicity gives
$A(1)-A(0)=\mathbf 1\{C\}$, where $C$ denotes a complier. Consequently
$\Delta_A=P(C)$ and
$\Delta_Y=E[(Y(1)-Y(0))\mathbf 1\{C\}]$; their ratio is
$E[Y(1)-Y(0)\mid C]$ when $P(C)>0$. The common limiting population
is the scientific premise. If eligibility itself changes the outcome
or the two sides approach different potential-outcome laws, this
decomposition does not identify the proposed action effect.

A finite authored construction exposes the exclusion issue without relying
on a fitted estimator. Let action probabilities below and above the threshold
be $.2$ and $.6$, and mean outcomes be USD~90 and USD~110. In one model,
the outcome mean is $80+50E[A\mid Z]$, so the causal coefficient is USD~50
and the reduced-form ratio is also 50. In another, the mean is
$84+30E[A\mid Z]+8Z$. It has exactly the same two action probabilities
and outcome means, but a causal coefficient of USD~30 and a direct
eligibility benefit of USD~8. With an expressly stipulated incremental
cost of USD~40, the implied value of the action is $+10$ or $-10$.
These equal summary moments are not asserted to be equal full distributions.
They demonstrate which additional restriction the ratio uses
\citep{exec13_ivrd_exact}.

More instruments do not automatically validate exclusion. Consider
independent fair signs $Z_1,Z_2,V,U$ and the structural equations
\[
 A=Z_1+2Z_2+V,\qquad
 Y=30A+10(Z_1+2Z_2)+U.
\]
An intervention changing $A$ while holding its causes fixed has coefficient
30. Nevertheless $Y-40A=U-10V$, so both population instrument moments
$E[Z_j(Y-40A)]$ equal zero. Instrument direct effects proportional to
their first-stage contributions escape this population overidentification
comparison. The exact sixteen-state enumeration checks the calculation;
it is not a vendor test run or a claim about a particular finite-sample
rejection probability. First-stage relevance, overidentifying restrictions
and causal instrument validity therefore need separate fields and explanations
\citep{exec13_ivrd_exact}.

\subsection{Stata already provides weak-instrument sets}

The inspected \texttt{ivregress} manual documents two-stage least squares,
limited-information maximum likelihood and generalized method of moments.
The postestimation manual distinguishes endogeneity tests, first-stage
statistics, overidentification tests and \texttt{estat weakrobust}.
For just-identified models the latter defaults to an Anderson--Rubin test;
for overidentified models it defaults to a conditional likelihood-ratio
test, with the documented generalization for a nonhomoskedastic covariance
estimate. The requested null concerns coefficients of the endogenous
regressors. Confidence-set output through \texttt{ci} is restricted to
a single endogenous regressor. A multiple-coefficient hypothesis test
must not be exported as though it were a fully computed joint confidence
region~\citep{exec13_ivrd_stata}.

The manual explicitly describes five confidence-set forms: a finite
interval, a union of finite intervals, a union that can include infinite
intervals, the whole real line and the empty set. Thus preserving
disconnected or unbounded answers is already part of a documented
commercial alternative. The proposed library's general quadratic-set
operation does not by itself establish a statistical advance over these
procedures. Its potential contribution is a common scientific contract
across tools, complete request accounting and a decision layer that
preserves the uncertainty and its construction.

The numerical distinction matters. For some cases Stata computes the
inversion directly; for nonhomoskedastic overidentified CLR inference it
uses a grid. The inspected documentation specifies an initial grid of
500 points, default bounds five conventional standard errors around the
estimated coefficient, adaptive expansion, a default maximum of 20
iterations, and options controlling bounds and adaptation. Where critical
values require simulation, the default is 25,000 draws and a seed can be
supplied. These are documented controls, not measured execution costs.
An adapter should retain the requested and final bounds, grid size,
simulation seed/state, replications, covariance specification and all
interval components~\citep{exec13_ivrd_stata}.

The manual distinguishes a model-rejecting empty Anderson--Rubin set
in an overidentified model from an apparent empty answer produced by
insufficient grid bounds in other cases. This prevents a downstream
agent from reading every empty set as scientific rejection. A simple
authored rectangle makes the danger concrete: $M\in[1,2]$ and
$Q\in[-1/2,1/2]$ imply the ratio relation set
$(-\infty,-2]\cup[2,\infty)$. A grid confined to $[-1,1]$ finds no
accepted points despite a nonempty set. This is an exact geometry example,
not a reproduction of Stata's AR or CLR algorithm. It motivates recording
whether the numerical search actually resolved the declared domain
\citep{exec13_ivrd_exact}.

The official feature pages associate weak-instrument inference with Stata~19
and identify its earlier StataNow availability in the Stata~18 feature
listing. That evidence does not imply that every fixed Stata~18 installation
contains the same command or that the mutable manual is identical across
updates. The deployment manifest must bind the actual installed edition,
update and command help. The inspected licensing page distinguishes an
identified single user from a specified number of concurrent network users.
It does not settle unattended service accounts, embedding, redistribution
or the commercial terms of a proposed MCP service. Those rights and prices
need a relevant license review before a licensed integration trial
\citep{exec13_ivrd_stata,exec13_ivrd_license}.

\subsection{Systems, diagnostics and model-implied instruments are different alternatives}

SAS/ETS~14.1's \texttt{PROC SYSLIN} documentation supplies a useful
systems-oriented comparison. It documents 2SLS, LIML, K-class, MELO,
SUR, 3SLS and full-information alternatives. Its technical discussion
states the instrument correlation and orthogonality requirements. The
\texttt{INSTRUMENTS} statement supplies all listed instruments to each
first-stage regression; an endogenous variable used as a regressor is
declared in \texttt{ENDOGENOUS}. The worked 2SLS example explains the
\texttt{FIRST} option and warns that its ANOVA and $R^2$ need a different
interpretation from ordinary least squares. A consumer therefore needs
the equation system, restrictions, instrument lists and actual estimation
sample, not just a coefficient table. These inspected pages establish
documented operations; they do not settle weak-IV set inversion or RD
support in every SAS edition. The incomplete download of a later manual
is excluded from technical evidence~\citep{exec13_ivrd_sas}.

EViews' instrument-diagnostic documentation distinguishes instrument
summary, subset orthogonality, regressor endogeneity, weak-instrument
diagnostics and a GMM breakpoint test. The summary can expose instruments
dropped for collinearity. Its weak-instrument view reports Cragg--Donald
statistics, available Stock--Yogo critical values and moment-selection
criteria for the stated estimation methods. The manual itself qualifies
the Cragg--Donald statistic as applicable to TSLS and other K-class
estimators while also displaying it for GMM comparison. The proposed
interface should carry that diagnostic qualification. A relevance screen
is neither a confidence set robust to weak identification nor verification
of exclusion. The extracted document's equations contain image elements;
only inspected textual definitions and scope are used here, without
reconstructing unread formulas~\citep{exec13_ivrd_eviews}.

IBM SPSS Statistics~30 documents a 2SLS procedure requiring the Standard
edition or Regression option. It exposes explanatory and instrumental
variable lists, coefficient/covariance output and command syntax. A
regressor omitted from the instrument list is treated as endogenous;
the documentation normally includes the exogenous explanatory variables
in both lists. This is useful for a guided analysis interface, but a
generic menu explanation calling a fitted proxy optimal cannot replace
the identifying assumptions, rank conditions or the appropriate covariance
argument. The inspected procedure page does not establish all-version
weak-IV or RD coverage. Its listed model-based intervals must retain
their own assumptions when exported~\citep{exec13_ivrd_spss}.

RATS~11 provides a programmable alternative. Section~2.5 of its 2025
user's guide specifies the \texttt{INSTRUMENTS} list and \texttt{INST}
option across its regression procedures, with system operations using
the documented GMM route. It writes the orthogonality condition explicitly
and distinguishes the covariance choice through \texttt{ROBUSTERRORS}
and other options. No variable, including the constant, is automatically
included in the instrument list. This small semantic difference can
change a translated model. A generated command file must therefore
declare the constant, included exogenous variables, excluded instruments,
sample window and covariance options explicitly. The source does not
support a claim that an arbitrary RATS program is a valid IV design,
or that every inference extension has been inspected~\citep{exec13_ivrd_rats}.

JMP~19's September 2025 multivariate manual documents model-implied
instrumental-variable 2SLS inside its structural-equation modeling
platform. It estimates equations separately and reports equation-level
Sargan diagnostics and the instrument set for each equation. This is a
valuable alternative when the analyst begins with a supplied structural
model. The instruments are implied by that model; discovering a valid
credit encouragement from an unrestricted business brief is a different
task. The documented local-misspecification advantages are not a theorem
that every structural misspecification or invalid instrument is harmless.
Edition entitlement and an actual scripted export remain to be checked
for the intended installation. A model-implied instrument should be
identified as such in the shared contract~\citep{exec13_ivrd_jmp}.

MATLAB's \texttt{ivx}/\texttt{iv4}, PcGive's instrument-estimation path,
and LIMDEP/NLOGIT remain named acquisition leads. Their essential technical
pages were inaccessible through the bounded official routes used here.
Consequently this comparison makes no technical equivalence, absence,
performance or deployment-rights claim for them. The concrete follow-up
is an edition-matched original manual and license record, followed by
the same instrument, covariance, numerical-set and export comparison.
This omission limits commercial coverage; it does not invalidate the
supported findings for the acquired sources~\citep{exec13_ivrd_survey}.

\subsection{An open RD package on a commercial host has two reuse decisions}

The pinned upstream Stata \texttt{rdrobust} help and source identify
version~11.1.1, dated 1 October 2026. They document local-polynomial
point estimation, robust bias-corrected inference, sharp and fuzzy
designs, distinct estimation and bias bandwidths, covariates and several
variance options. The help distinguishes conventional inference,
bias-corrected estimates with conventional variance, and bias-corrected
estimates with robust variance. A result containing only one coefficient
and standard error loses that distinction. Its MSE and coverage-error
bandwidth choices also need their labels retained. The inspected source
requires Stata version~16 semantics; this is a declared source requirement,
not a completed compatibility trial~\citep{exec13_ivrd_rdrobust}.

The \texttt{masspoints} options inspect repeated running-variable values.
The default adjustment enlarges preliminary bandwidths to include a
minimum number of distinct values, with ten as the documented default
and an explicit override. The source computes unique-score counts on
both sides for manual as well as automatically chosen bandwidths.
Its fuzzy branch rejects the fully degenerate constant treatment case
with no jump. That exact guard does not by itself establish coverage
when the first stage is merely small. Nor is robust bias correction
synonymous with weak-IV robustness. A matched trial must retain both
bandwidth uncertainty and weak local treatment uptake in its design
\citep{exec13_ivrd_rdrobust}.

The companion pinned \texttt{rddensity} help describes local-polynomial
density manipulation tests, bias-corrected and conventional statistics,
mass-point adjustment and complementary binomial testing. Its graphical
path requires \texttt{lpdensity}. Those are useful diagnostics, not a
certificate that applicants could not manipulate scores or that all
potential outcomes are continuous. A nondetectable density change can
coexist with another discontinuous policy; a density rejection can
reflect a measurement or design issue that needs investigation. The
proposal should retain test statistic, bandwidth, cutoff, support counts,
test family and qualification instead of converting a $p$-value into
a Boolean ``causal design verified'' field~\citep{exec13_ivrd_rddensity}.

Discrete scores show the remaining information problem. Suppose the
observed support is $\{-2,-1,1,2\}$, eligibility starts at zero, and
observed means are $80,90,130,140$. Let the treated response function
be $120+10s$. Two smooth untreated functions are
\[
 \mu_{0,\delta}(s)=100+10s+
             \frac{\delta}{2}(s+1)(s+2),\qquad \delta\in\{0,40\}.
\]
Both agree with the observed untreated values at $-2$ and $-1$.
Under sharp assignment the treated observations agree as well. Yet
the causal discontinuity at the unobserved cutoff is $20-\delta$,
which is $+20$ or $-20$. Repeating these same scores indefinitely
does not supply observations tending to zero. Smoothness without a
quantitative shape restriction does not choose between these laws.
A justified curvature bound, prespecified response model, local
randomization law or new score resolution is additional information.
Counting mass points is useful; it cannot create the missing support
\citep{exec13_ivrd_exact}.

The two upstream license files declare GPL version~3. The proposed
default is a reviewed adapter to a pinned distribution, with the
applicable license obligations assessed for the actual distribution
and use. Stata's host license is a separate issue. A port would need
to preserve the polynomial basis, kernel, bandwidth selection, bias
variance, cluster corrections, support accounting and missing-data
rules, with independent reference comparisons and an appropriate
license decision. Rewriting a few coefficient formulas would not
reproduce the package. Direct adoption inside a licensed Stata
workflow may be preferable for an existing user; a native R/Python
path may be preferable for another consumer, but cross-language
equivalence is an empirical obligation~\citep{exec13_ivrd_rdrobust,exec13_ivrd_rddensity}.

\subsection{The proposed product must compare complete decisions}

A useful IV/RD request carries the business decision, target population,
action and outcome units, decision and outcome clocks, cutoff rule,
running-variable precision, instrument origin, exclusion and compliance
assumptions, estimation sample and covariance design. The result retains
the target actually estimated, first-stage and support diagnostics,
every interval component, numerical search status, all failed requests
and the source/engine identities. A report must distinguish a
nonidentified effect, an unsupported operation, a scientifically
refused request and an unfinished numerical search.

The current local reference service can preserve generic supplied
moment and sensitivity sets, but it does not implement the acquired
commercial procedures or the pinned RD packages. The proposed
integration should begin with one licensed IV path and one pinned
RD adapter, each compared with a separate exact or small numerical
reference. The trial should include strong, weak and zero first stages;
valid and invalid exclusion; correlated and heteroskedastic observations;
coarse scores; manipulated cutoffs; alternative bandwidths; empty,
disconnected and whole-line sets; and grid exhaustion. It should
record all model fitting, conversion, failed-request and review costs.
These are proposed comparison cases, not completed vendor experiments.

The investment decision is whether those paths reduce errors and
repeated manual work enough to justify the integration. It is not
established by a longer list of estimator names. Stata's existing
weak-IV sets are a substantive baseline; SAS, RATS, EViews, SPSS
and JMP supply different supported entry points. The near-term
opportunity is to preserve their scientific meaning, expose their
limits, and connect a locally justified effect to a decision with
explicit costs. Measured savings and bank returns remain later
empirical questions~\citep{exec13_ivrd_survey,exec13_ivrd_exact}.
\label{exec13:commercial-iv-rd-end}


% SOURCE docs/release/2026-10-29/source/docs/writing-repair/2026-10-29/src/87ae10133a61865d.tex CHAR 0
\section{Commercial survival software and repeated banking interventions}
\label{exec15:commercial-survival}

A servicing team observes default, prepayment, hardship contact and changing
arrears over a loan's life. It asks whether repeated assistance will improve
customer outcomes and the value of the portfolio. Several familiar software
operations sound relevant: a survival curve, a Cox regression with changing
covariates, an average treatment effect on survival time, a competing-event
model and a latent growth model. They do not answer the same question.
An adequate comparison begins with the decision, eligible population, common
time origin, event definitions and intervention history. Only then can the
proposed library decide which commercial operation is a valid comparator.

This bounded review develops eight families from official technical text:
Stata, SAS/STAT, IBM SPSS, JMP, Minitab, NCSS, Statistica and Mplus. MATLAB
remains a ninth candidate because both attempted official method pages were
unavailable. Several sources identify an edition; the exact NCSS and Mplus
editions and the Minitab update remain unresolved. A mutable manual acquired
after 7 October 2026 does not establish its exact content at that literature
cutoff. Consequently the comparison supports its inspected method claims,
not an exhaustive historical market census or installed-product qualification.
No commercial engine was executed, and no price, deployment entitlement or
performance advantage is inferred from documentation
\citep{exec14_survival_survey}.

\subsection{Time, events and causal targets belong in the interface}

Let $T^g$ be the first terminal-event time under a declared feasible policy
$g$, and $J^g$ the event type, such as default or prepayment. A first-event
survival probability $P(T^g>t)$ differs from the default cumulative incidence
$F_D^g(t)=P(T^g\leq t,J^g=D)$. The restricted mean time without either event
is $E[\min(T^g,\tau)]$, for a specified horizon $\tau$. It is neither the
unrestricted mean default time nor a ratio of conditional hazards. A policy
may change prepayment as well as default, so changing the target can change
the apparent benefit. These objects are proposed contract fields, not
automatically identified quantities.

For a discrete two-period construction, write $d_k^g$ and $p_k^g$ for the
conditional probabilities of default and prepayment during period $k$ among
loans still active at its start. They are mutually exclusive, with
$d_k^g+p_k^g\leq1$. Set $S_0^g=1$. Direct conditioning gives
\[
 S_k^g=S_{k-1}^g(1-d_k^g-p_k^g),\qquad
 F_D^g(K)=\sum_{k=1}^K S_{k-1}^g d_k^g.
\]
The analogous sum with $p_k^g$ gives prepayment incidence. The recursion is
an accounting identity within the stipulated policy law. Identification of
that law from bank observations requires additional causal and observation
assumptions. In particular, replacing the recursion by
$1-\prod_k(1-d_k^g)$ removes prepayment while retaining the same default
hazards. That is a different construction; the unchanged-hazard premise
does not follow from observing the original competing-event process
\citep{exec14_survival_exact}.

Suppose both periods have default probability $1/10$ under two generated
policies. Policy A has prepayment probability $2/5$, while B has probability
$1/10$. After two periods A has default incidence $.15$, prepayment incidence
$.60$ and active probability $.25$. B has corresponding probabilities $.18$,
$.18$ and $.64$. Removing prepayment and retaining the default hazard gives
$.19$ under either policy. All three default numbers answer different
questions. A lower observed default incidence under A can arise because more
loans leave through prepayment, even though the conditional default hazard
is identical. Full enumeration of both period outcomes agrees with the
recursion \citep{exec14_survival_exact}.

Value need not rank the policies by default incidence. Stipulate USD~10 per
active loan at the start of each period, USD~100 lost at default and USD~20
lost at prepayment, with no discounting. Expected active periods are $1.5$
and $1.8$, so values are
\[
 V_A=10(1.5)-100(.15)-20(.60)=-12,\qquad
 V_B=10(1.8)-100(.18)-20(.18)=-3.6.
\]
The policy with fewer defaults has lower value under these exposed prices.
These are authored scenarios, not estimated bank effects or claims that
prepayment always causes a loss. The interface should return the joint event
accounting and the declared value map, allowing different accepted prices
without changing the scientific target \citep{exec14_survival_exact}.

\subsection{Stata offers causal survival comparators with explicit limits}

The inspected \texttt{stteffects} manuals document ATEs, effects on the treated
and potential-outcome mean survival times through regression adjustment,
inverse-probability weighting and combined methods. These are substantial
existing causal-estimation capabilities. Comparing the proposed library only
with an unadjusted survival plot would understate the commercial alternative.
The manuals also expressly exclude delayed entry, time-varying covariates and
multiple-record data from this command family. A documented baseline-treatment
survival estimator therefore cannot be silently relabeled as an estimator of
a repeated collections policy \citep{exec14_survival_stata}.

Three scientific conditions remain distinct: conditional independence of
treatment and potential outcomes given the supplied covariates, treatment
overlap and correct adjustment for censoring. The documentation distinguishes
likelihood-adjusted censoring in LAC-IPWRA from weighted-adjusted censoring
in WAC-IPWRA. In the inspected implementation, weighted censoring requires
random, correctly modeled censoring and a censoring process that does not
vary by treatment level. The choice is consequential for a servicing example
where treatment changes the chance that records disappear. Calling an
estimator doubly robust does not remove an identification premise or extend
the documented censoring regime \citep{exec14_survival_stata}.

The proposed adapter should retain the chosen estimand, treatment/control
labels, outcome distribution, treatment model and censoring adjustment as
separate fields. It should export propensity scores, censor-survival
probabilities, conditional means, covariance type and the actual estimation
sample. The manual documents an overlap tolerance that can stop estimation
and an option identifying offending observations. It also distinguishes
robust, clustered, bootstrap and jackknife uncertainty. A refusal at an
overlap boundary belongs in the result, not in an omitted request. Balance
and overlap diagnostics examine supplied observed covariates; they do not
authenticate an unmeasured-confounding assumption
\citep{exec14_survival_stata}.

Stata's separate \texttt{stcrreg} command models the subhazard of a selected
event with specified competing events. It permits single or multiple records
but not multiple failures per subject. Its event codes, time interactions,
cluster settings and cumulative-incidence interpretation require a different
conversion contract from \texttt{stteffects}. Recurrent missed payments and
terminal default are different event structures even when the raw table has
the same number of rows. The first proposed commercial trial should use a
supported fixed-treatment case and a separately specified competing-event
case; it should not claim that either demonstrates arbitrary dynamic-policy
evaluation \citep{exec14_survival_stata_competing}.

\subsection{SAS and JMP illustrate why competing-event outputs differ}

SAS/STAT 14.1 PHREG documents counting-process input in intervals $(T_1,T_2]$,
with covariates held fixed within each record and event status at the end.
Its examples cover repeated events and delayed entry. This is useful for
representing servicing records faithfully, provided subject identity, actual
entry time and covariate availability are preserved. The input representation
does not itself identify the effect of intervening on a history. It specifies
who belongs in a model's risk set \citep{exec14_survival_sas}.

The same edition documents a distinct subdistribution-hazard pathway through
\texttt{EVENTCODE}. Subjects with competing events remain in modified risk
sets with weights; they should not be exported as ordinary event-free subjects.
The manual specifies a sandwich covariance, an explicit cluster-ID route,
cumulative-incidence predictions and pointwise interval settings. It also
states that several options and statements are ignored in this mode, including
\texttt{ASSESS}, \texttt{BAYES} and \texttt{RANDOM}. A generic wrapper that
accepts every option and labels it effective would misdescribe the analysis.
The adapter must preserve the effective procedure configuration and distinguish
a pointwise band from a simultaneous or policy-value guarantee
\citep{exec14_survival_sas}.

JMP 18.1's inspected Competing Causes page describes cause-specific survival
estimation by censoring events from other causes and states independence for
its cause-wise interpretation. Its Omit Causes control recalculates the
display to illustrate removal of selected causes. In the banking construction,
removing prepayment and holding default hazards unchanged yields $.19$, while
actual default incidences are $.15$ and $.18$. Thus a cause-removal display
is not interchangeable with the observed cumulative default probability.
Nor does the control establish that a bank can remove prepayment without
changing default mechanisms. This is a scoped comparison of the documented
platform, not an assertion that every JMP survival capability shares that
interpretation \citep{exec14_survival_jmp,exec14_survival_exact}.

The proposal should therefore avoid importing a column merely named
``survival'' or ``risk.'' Its descriptor needs the event set, horizon,
conditional population, risk-set convention, cause-removal premise where
used, and uncertainty definition. SAS's subdistribution output and the
inspected JMP cause-wise display are useful alternatives for different
questions. A matched experiment first harmonizes the question and only then
compares numerical estimates, diagnostics and workflow cost.

\subsection{Time-dependent regression must preserve the decision clock}

IBM SPSS Statistics 30 documents the system time variable \texttt{T\_},
both for explicit time interactions and for segmented measured covariates.
Its weekly example combines logical interval indicators with measurements.
Statistica 14.2 separately offers survival-time inputs, explicit time-function
covariates and counting-process start/stop inputs. These are useful ways to
represent changing records, but a change of input form is not evidence of a
causal dynamic-regime estimator. Measurement time, availability time and the
action time must still be checked independently
\citep{exec14_survival_ibm,exec14_survival_statistica}.

Consider 100 stipulated loans at a common initial eligibility date. Twenty
default before day two. Among the 80 survivors, 40 receive assistance at day
two and 40 do not; ten in each group subsequently default. Labeling exposure
from the initial date by whether assistance ever occurs gives risks $10/40$
and $30/60$, an apparent difference of minus 25 percentage points. The
future-treated category necessarily excludes the twenty early defaults.
A day-two landmark comparison gives $10/40$ in each arm, but now concerns
the survivor population. Neither rewriting the label nor dropping early events
recovers a baseline assignment-policy effect. The construction diagnoses a
clock and denominator error; it is not an estimate of Cox-regression bias
\citep{exec14_survival_exact}.

A converter should split records at declared changes while preserving which
covariate value was known immediately before the decision. In a $(s,t]$
record, the end event and any new information at $t$ must follow an explicit
ordering convention. A measurement collected after an event cannot be carried
back into its risk interval. A policy scheduled for day two needs a definition
for loans that default earlier. These are scientific design choices, not
formatting repairs; the software should reject ambiguity rather than choose
an ordering from the row sort.

Minitab's inspected counting-process method table makes several conversion
checks concrete. It distinguishes fixed predictors from time-dependent
repeated measures, records categorical coding and tie adjustment, distinguishes
robust from model covariance, and describes stratification and proportional hazards diagnostics. It also documents excluded rows with missing or
nonpositive response times and an error when entry exceeds exit. Consequently
the input row count is insufficient: a proposed adapter should reconcile
subjects, intervals, events and exclusions, and preserve the effective
covariance and tie method. This is documentation-based design; exact edition
and runtime parity still need qualification \citep{exec14_survival_minitab}.

\subsection{Sequential confounding is not removed by a time column}

An exact two-decision construction shows the additional causal work. Let
$A_0$ be an early assistance action randomized with probability $1/2$.
Let $L$ be later arrears, with probabilities $.8$ under $A_0=0$ and $.2$
under $A_0=1$. The later action $A_1$ has probability $.8$ when $L=1$ and
$.2$ when $L=0$. Define final default probability by
\[
 P(Y=1\mid L=l,A_1=a)=.1+.4l-.05a.
\]
Independent randomizers generate the stated conditional probabilities. This
fully specified law supplies consistency, sequential exchangeability and
positive treatment support for this example. None is inferred from a fitted
hazard model or a software diagnostic \citep{exec14_survival_exact}.

Marginally $P(L=1)=1/2$. Later assistance selects higher-risk histories:
$P(L=1\mid A_1=1)=.8$ and $P(L=1\mid A_1=0)=.2$. Hence observed default
risks are $.37$ with $A_1=1$ and $.18$ with $A_1=0$, despite the negative
conditional action effect. Under fixed policies, direct summation instead gives
\[
 \begin{split}
 P(Y^{0,0}=1)&=.2(.1)+.8(.5)=.42,\\
 P(Y^{1,1}=1)&=.8(.05)+.2(.45)=.13.
 \end{split}
\]
Their difference is $-.29$, while the observed late-action difference is
$+.19$. The two policy values are also reproduced by averaging the default indicator $Y$ over the complete
sixteen-state law with weights
\[
 \frac{\mathbf 1\{A_0=a_0,A_1=a_1\}}
 {P(A_0=a_0)P(A_1=a_1\mid L)}.
\]
This is a population identity for the authored law, not a fitted weighting
study or an asserted finite-sample coverage result
\citep{exec14_survival_exact}.

Conditioning on $L$ can remove confounding of the later action while also
holding fixed a variable changed by the earlier action. In this construction,
the controlled $A_0$ difference at fixed $(L,A_1)$ is zero, although the
two-regime contrast is not. A single coefficient from a regression containing
$A_0,L,A_1$ cannot be interpreted as the total effect of the repeated policy
without further calculation and assumptions. The proposed library needs an
explicit longitudinal target, history-indexed assignment/observation laws
and a valid standardization or other justified procedure. A counting-process
adapter can supply faithful data to such work; it is not a substitute for it.

Mplus's inspected Chapter 6 offers distinct discrete-time survival, Cox,
parametric baseline-hazard and latent-factor survival examples. In the
discrete example, post-event periods and dropout can share a missing-value
code. A bank adapter should retain their different reasons before conversion.
The latent-factor example makes numerical integration and its computational
burden explicit. Its measurement model is a substantive alternative when
supported, but fitting a latent factor does not authenticate causal confounding
control. These examples justify a narrow model-and-output comparator, while
the chapter's exact edition, intended runtime and rights remain unresolved
\citep{exec14_survival_mplus}.

\subsection{Censoring and failed uncertainty calculations remain outcomes}

A small horizon-observation example separates an outcome model from its
observed cases. Half the population is low risk with default probability $.1$;
half is high risk with probability $.6$. All low-risk outcomes and one quarter
of high-risk outcomes are observed. Observation is stipulated independent
of default conditional on this recorded stratum. Full-cohort risk is
$.5(.1)+.5(.6)=.35$, while complete-case risk is
\[
 \frac{.5(.1)+.5(.25)(.6)}{.5+.5(.25)}=.20.
\]
Inverse observation weighting restores $.35$ under the known probabilities.
It cannot do so by arithmetic alone if the stratum is unavailable, the
probabilities are wrong or observation depends on unrecorded future outcomes.
This finite missingness construction motivates the analogous checks for a
time-indexed censoring model; it does not establish a continuous-time estimator
\citep{exec14_survival_exact}.

NCSS's Cox manual supplies a useful reporting comparator even without a
qualified integration. Its run summary contains processed and estimation
rows, likelihood iterations, achieved convergence and completion status.
The described stopping statistic is the maximum relative coefficient change
on the last iteration. A finite number in a coefficient table is therefore
not enough to classify a run as converged. The proposal should preserve
these diagnostics and distinguish numerical nonconvergence, a scientifically
unidentified target and an unsupported method. The source is an undated
official chapter; it supports the inspected reporting semantics, not a
claim about a particular installed NCSS release \citep{exec14_survival_ncss}.

Uncertainty calculations need their own failure denominator. Suppose eight
subjects include only one with a required rare treatment history. An ordinary
eight-draw subject bootstrap omits that subject with probability
$(7/8)^8\simeq.3436$. An unregularized mean for that history is undefined in
those resamples. Discarding them and silently drawing replacements conditions
the replication distribution on containing the history. It is not the
unconditional bootstrap originally declared. This exact occupancy calculation
does not claim a 34.36\% failure rate for any vendor command: it specifies a
stress case for the proposed wrapper. A cluster bootstrap, stratified design,
regularization or a restricted target is a substantive alternative with its
own assumptions, not a license to hide failures \citep{exec14_survival_exact}.

\subsection{Build a target-preserving comparator before a product claim}

The six-operation reference revision examined here handled declared graph
syntax, a supplied-DAG backdoor case, supplied-moment inversion, deterministic
sensitivity sets and declared observation clocks. It did not expose fitted
survival, competing-risk or dynamic-regime estimation. The later servicing
study in Chapter~\ref{gap29:servicing} supplies bounded fitted research;
Chapter~\ref{gap29:public} demonstrates fitted-summary inference across
interfaces. Neither result retroactively changes this probe or establishes a
general raw-history survival service. The inspected
commercial procedures already supply serious estimation, diagnostics and
reporting alternatives; the proposed differentiator must be demonstrated in
target preservation, supported causal scope and transparent uncertainty,
rather than assumed from having an agent interface
\citep{exec14_survival_survey}.

The immediate route is a public/generated comparator contract. Begin with a
fixed-treatment survival target supported by the chosen procedure; add a
competing-event case with separately declared incidence and cause-removal
targets. Verify subject identity, time units, entry/exit boundaries, event
codes, treatment levels, selected rows and output units. Then inspect complete
failure trajectories, censoring diagnostics and uncertainty settings. Repeated
policies enter a separate experiment with an explicit sequential identification
argument and exact reference law. No unsupported request should be routed to
the nearest command by name.

Stata and SAS warrant consideration as licensed numerical comparators;
IBM, Statistica and Minitab supply useful clock, representation and diagnostic
comparisons; JMP clarifies cause-removal semantics; NCSS exposes convergence
reporting; and Mplus supplies survival/measurement-model alternatives.
These are proposed uses of documented paths. Proprietary code should not be
ported under an assumed license. A future adapter or direct-use decision
requires the selected edition, authorized execution mode, permitted exports
and applicable redistribution and benchmark-publication terms. MATLAB and
uninspected extensions stay unresolved. Missing evidence is not a negative
feature finding \citep{exec14_survival_survey}.

Only after the public/generated prototype preserves these targets and makes
the intended comparison useful should a paid runtime study or authentic bank
experiment be commissioned. Its cost includes the license, conversion,
fitting, uncertainty replication, failed attempts, review and maintenance.
No current quotation or bank ROI is available from this review. That is a
specific later input, while the present proposal can already specify the
scientific comparisons, acceptance cases and unsupported uses.
\label{exec15:commercial-survival-end}


% SOURCE docs/release/2026-10-29/source/docs/writing-repair/2026-10-29/src/22fb3e3839fcd95c.tex CHAR 0
\section{Commercial survival options against the proposed service}
\label{exec16:commercial-matrix}

The preceding survival analysis compares several procedures on their own
terms. A procurement decision needs a second comparison: what each inspected
edition documents, what the tested local reference revision did, and what
the proposal still asks an implementation to supply. These are three different
columns. In particular, an operation's appearance in a product proposal is
not evidence that its reference service computes it. The tested local
reference revision implemented a bounded supplied-moment residual inversion,
a supplied DAG back-door construction for independent units, graph syntax,
deterministic sensitivity and observation-clock checks. Its dispatcher returned
\texttt{unsupported} for the three authored fixed-treatment survival,
competing-incidence and sequential-policy identifiers. This observation is a
test of the local dispatcher, not a vendor execution or an attempted encoding
of every survival target. That revision's request enum also had no separately
named survival operation. These facts establish a gap for those tested cases.
Later bounded research and fitted-summary interface results do not qualify a
general raw-history survival estimator; they must nonetheless remain visible
when describing the proposal's present evidence.
\citep{exec16_commercial_runtime}

For a fixed treatment at a common origin, Stata's inspected
\texttt{stteffects} documentation already describes average effects, effects
on the treated and potential mean survival outcomes with regression,
weighting and combined estimators. It documents overlap and censoring
diagnostics. This is a substantive commercial comparator. The same command
family expressly excludes delayed entry, time-varying covariates and
multiple-record data; its separate \texttt{stcrreg} subhazard path has another
target and conversion contract. The proposal should first compare a
documented fixed-treatment profile and a separate competing-event profile,
retaining treatment labels, analysis sample, censoring adjustment and
uncertainty. A dynamic servicing-policy claim would require additional
identification and implementation evidence. The inspected mutable manuals
do not verify a particular installed Stata update at the literature cutoff.
\citep{exec14_survival_stata,exec14_survival_stata_competing}

SAS/STAT 14.1 PHREG documents interval records, time-dependent covariates,
delayed entry and recurrent-event examples, as well as an
\texttt{EVENTCODE} subdistribution path with modified competing-event risk
sets. It is a useful licensed comparator for faithful record preparation and
cumulative-incidence output. Some statements and options are ignored in the
documented subdistribution mode; a bridge must report effective settings,
not only submitted text. The inspected SAS route does not by itself identify
an intervention on a repeated assistance policy. The much larger later PHREG
manual was not acquired within the fixed transfer bound, so later-edition
equivalence remains unreviewed rather than assumed.
\citep{exec14_survival_sas}

Other inspected products can answer narrower conversion questions. IBM SPSS
Statistics 30 documents \texttt{T\_} time expressions and measured-covariate
segments. Statistica 14.2 distinguishes survival-time, time-function and
start--stop forms. Minitab's counting-process help records tie adjustment,
covariance choice, excluded rows and entry/exit errors. These are useful for
testing the proposed clock and row-reconciliation layer, but a time-varying
column is not a sequential policy estimator. JMP 18.1's inspected
Competing Causes display treats other events as censored in cause-wise
analysis under its stated independence premise; \emph{Omit Causes}
recalculates a removal illustration. It is not the same object as an
observed default cumulative incidence or an established intervention on
prepayment. A comparison that reads either output as a policy value has
changed the scientific question before considering software quality.
\citep{exec14_survival_ibm,exec14_survival_statistica,exec14_survival_minitab,exec14_survival_jmp}

NCSS's inspected Cox chapter contributes a reporting comparator: processed
and used rows, iterations, achieved convergence and completion status. A
finite coefficient without convergence is not a completed scientific result.
Mplus's inspected chapter gives discrete-time, Cox, parametric and latent-factor
survival examples. It illustrates both a richer modeling alternative and the
need to separate dropout from post-event missingness; a latent measurement
model does not verify the absence of hidden confounding. Exact NCSS and Mplus
editions are unresolved. MATLAB remains a candidate, but the two attempted
official method originals returned HTTP 403. No feature-absence or quality
ranking follows from that acquisition outcome.
\citep{exec14_survival_ncss,exec14_survival_mplus,exec16_commercial_matrix}

The comparison therefore does not count a vendor's documented Cox feature
against a proposed causal-policy feature as though both were exchangeable.
For each profile, the proposed library must bind treatment and event
definitions, time origin, eligible subjects, history and availability clock,
observation/censoring law, estimand, uncertainty and value units. The initial
public/generated trial can compare fixed-treatment survival where a chosen
licensed method actually applies, then a competing-event example with separate
incidence and cause-removal outputs. A generated start--stop fixture tests
identity, excluded rows, ties and event ordering. A sequential-policy example
requires its own assignment and observation law. Admission of a selected
estimator still requires implementation and qualification for that declared
profile.
The earlier two-period competing-event and sixteen-state sequential examples
provide exact references, not vendor performance data.
\citep{exec14_survival_exact,exec16_commercial_matrix}

This route also determines what evidence a paid evaluation would buy. The
prototype should first show that the selected target and conversions matter
on public or generated inputs. Then an owner could select an exact commercial
edition, obtain the intended entitlement and confirm automation, export,
embedding and benchmark-publication rights. A frozen matched run would compare
the same profile, input population, failure denominator, diagnostics,
uncertainty and full workflow cost. Prices and those rights are unavailable
in the inspected record, so the present comparison cannot choose a purchase
or claim a cheaper product. Missing historical edition bytes, uninspected
families and unavailable MATLAB originals keep overall coverage inconclusive.
They do not prevent the bounded prototype contract from advancing.
\citep{exec16_commercial_matrix}
\label{exec16:commercial-matrix-end}


% SOURCE docs/release/2026-10-29/source/docs/writing-repair/2026-10-29/src/0a2846d650819f7b.tex CHAR 0
\section{Commercial policy evaluation: what an offline score permits}
\label{exec16:commercial-policy}

A servicing manager comparing a contact policy with the existing workflow
needs a value for a specified action rule, population and horizon. Several
commercial tools expose quantities called policy value, reward, uplift or
expected utility. These quantities support different decisions. A static
treatment-assignment rule, a learning procedure that changes after each event,
a sequence of servicing actions, and an optimum within a supplied decision
model are distinct objects. The proposal should compare them at those
boundaries before deciding whether to buy a license, integrate an engine or
build a missing component.

The documented landscape already contains useful capabilities. Stata's
\texttt{cate} postestimation evaluates prespecified treatment-assignment
policies. The historical Azure Personalizer service documents offline
comparisons of learning policies. GeNIe and its SMILE engine provide decision
analysis and expected utilities in supplied influence diagrams. Amazon
Personalize documents held-out recommendation metrics. JMP's uplift platform
provides a separate targeting alternative. These are substantive comparators,
but they do not jointly establish a generally valid sequential causal-policy
service. The following comparison uses inspected official technical text;
none of these vendor engines was executed for this review.
\citep{exec13_ch_stata,exec16_cp_azure,exec16_cp_bayesfusion,exec16_cp_aws,exec13_ch_jmp_uplift}

\subsection{Four policy questions and one declared information clock}

For one binary treatment, let $X$ denote information available before treatment,
$Y(a)$ the potential outcome, and $\pi(x)\in[0,1]$ a prespecified probability
of assigning treatment. The static value is
\[
 V(\pi)=E\{\pi(X)Y(1)+[1-\pi(X)]Y(0)\}.
\]
This is the target displayed in Stata's policy-evaluation methods. The
command also compares two policies and accepts probability-valued policy
variables. Its synonym \texttt{estat tassigneval} should not be counted as
a separate method. The premises of the underlying treatment-effect analysis
still apply. A rule chosen by maximizing this same evaluation sample is a
different inferential target from a prespecified rule; a changing sequence of
actions is not represented merely by placing its last action in the treatment
column. \citep{exec13_ch_stata}

For sequential servicing, let $H_t$ be the observed history immediately before
action $A_t$, $R_t$ the subsequent net cash flow, $b_t(a\mid h)$ the logging
probability, and $\pi_t(a\mid h)$ the proposed policy. A finite-horizon target is
\[
 V_T(\pi)=E_\pi\left[\sum_{t=0}^{T-1}\gamma^tR_t\right],\qquad
 W_t=\prod_{s=0}^t\frac{\pi_s(A_s\mid H_s)}{b_s(A_s\mid H_s)}.
\]
The change-of-law argument needs consistent actions and outcomes, sufficient
histories for sequential assignment, positive logging probability wherever
the target policy acts, and an appropriate initial population and observation
law. Under a correctly specified observed-state experiment, the likelihood
ratios change the action distribution while leaving the shared environmental
mechanisms unchanged. Jiang and Li develop trajectory and stepwise importance
sampling and a sequential doubly robust construction; their discussion also
requires separating a data-derived target policy from the sample used for
evaluation. That result does not grant causal validity to arbitrary hidden
servicing states or misrecorded logging probabilities. \citep{jiang2016}

A supplied-model decision engine answers a third question: what action has
the highest expected utility if its probability tables, information sets and
utilities are correct? A recommendation metric answers a fourth: how well
does a learned ranking reproduce or predict held-out interactions? Either
can be valuable without identifying the effect of deploying a new servicing
policy. The commercial feature matrix should retain the target, source of
mechanisms, sampling unit and uncertainty claim for each quantity.

\subsection{A final-action adjustment can reverse a servicing decision}

Consider a generated two-stage law with $\gamma=1$. The first action $A_0$ is a fair binary
draw; the next observed state is $L_1=A_0$. The logger chooses $A_1=1$ with
probability $9/10$ after $A_0=0$ and $1/10$ after $A_0=1$. Define a terminal
net reward, in stipulated dollars, by
\[
 R=20A_0+100A_1-120A_0A_1-5(A_0+A_1).
\]
This law deliberately models an adverse interaction between repeated actions;
it is not an estimated banking mechanism. Its four paths are
\[
 \begin{array}{c|rrrr}
 (A_0,A_1)&(0,0)&(0,1)&(1,0)&(1,1)\\\hline
 P_b&1/20&9/20&9/20&1/20\\
 R&0&95&15&-10
 \end{array}
\]
For the policy setting both actions to one, direct substitution gives
$V=-10$. The full trajectory weight is twenty on path $(1,1)$ and zero
otherwise, hence $E_b[WR]=(1/20)(20)(-10)=-10$.
\citep{exec16_cp_exact}

Using only the history-conditional last-action ratio answers another question:
\[
 E_b\left[\frac{\mathbf1\{A_1=1\}}{b_1(1\mid A_0)}R\right]
 =\tfrac12(95)+\tfrac12(-10)=42.5.
\]
It retains the logger's first-action distribution. Using the marginal
last-action probability $P_b(A_1=1)=1/2$ instead gives
$2\{(9/20)95+(1/20)(-10)\}=84.5$. Neither positive number evaluates the
two-action target whose value is negative. The example specifies exactly
what a static-policy comparator can and cannot replace. More observations
of the same law would estimate the wrong target more precisely.
\citep{exec16_cp_exact}

There is also a precision cost even when the weighting is correct. Here
$E_bW=1$ and $E_bW^2=20$, so the population ratio
$(E_bW)^2/E_bW^2$ is $1/20$. Only one twentieth of trajectories follow
the target path in expectation. This ratio explains a loss of information;
it is not a finite-sample coverage guarantee or a universal effective sample
size. The number of logged events alone cannot settle whether a policy
comparison is precise enough for an investment decision.
\citep{exec16_cp_exact}

\subsection{Azure Personalizer is a historical workflow reference}

The acquired Microsoft pages describe Rank requests containing available
actions, action features, context and an event identifier, followed by a Reward
call. The inspected overview uses a reward between zero and one and permits
delayed feedback. It describes learning across users and expressly says that
the service does not maintain individual user histories. These are useful
workflow features, but a sequence of event-level calls is not itself a
specified customer-trajectory causal model. A monetary servicing objective
also needs an explicit reward transformation, treatment costs and horizon;
a click score cannot be silently substituted for discounted net recovery.
\citep{exec16_cp_azure}

The offline-evaluation pages compare the current online learning policy,
the application's default action, a random-action policy, uploaded alternatives
and policies found by optimization discovery. They display estimated average
rewards and confidence intervals, allow a policy to be downloaded or applied,
and recommend representative traffic and at least 50,000 events as a practical
starting point. The conceptual description replays chronological events with
a blank model for each learning policy and retains observed rewards when the
chosen action matches. A learning-policy configuration in this description
controls a changing learner; it is not necessarily a single fixed action rule.
\citep{exec16_cp_azure}

That conceptual account is insufficient to reconstruct a general estimator
or its uncertainty theorem. For an exact illustration, suppose two contexts
are equally common, a target always chooses action one, and its mean rewards
are 10 and 100. Let the logger's action-one probabilities be $0.9$ and $0.1$.
The target value is 55, but the mean among matching logged actions is
\[
 \frac{0.9(10)+0.1(100)}{0.9+0.1}=19.
\]
Weighting by the correct context-specific assignment probabilities recovers
55 in this population. This is a counterexample to matching alone under
nonuniform logging, not a claim that Microsoft's actual implementation uses
that incomplete estimator. The inspected pages do not supply all the
algorithmic and confidence-bound details needed to make such a runtime
claim. A successor comparison must acquire those details and verify the
logging distribution, retained-event denominator, replay updates and
policy-selection rule. \citep{exec16_cp_azure,exec16_cp_exact}

Availability changes the acquisition recommendation. The official pages state
that new resources became unavailable on 20 September 2023 and that the service
would retire on 25 August 2026. That retirement date precedes the proposal's
7 October 2026 cutoff. The pages recommend the open-source
\texttt{microsoft/learning-loop} project. Thus Personalizer is a historical
workflow comparator rather than a proposed new hosted dependency. The named
successor needs its own pinned code, rights and estimator review; the migration
link does not establish feature or statistical parity. The pages were acquired
after the cutoff and are labeled as archived documentation, with their
displayed update dates and source metadata retained. \citep{exec16_cp_azure}

\subsection{Recommendation quality, incremental value and capacity}

Amazon Personalize's inspected technical page defines train/test constructions,
metrics such as precision, recall and normalized discounted cumulative gain,
and a metric named \texttt{average\_rewards\_at\_k}. The latter normalizes
rewards from interactions with selected recommendations relative to rewards
from interactions with recommendations. Its name does not make it an estimate
of the population intervention contrast. The page also warns against comparing
solution metrics across different training data, and provides CLI and SDK
routes to retrieve metrics. Those are meaningful operational capabilities to
compare, while their documented target remains distinct from a causal policy
effect. This mutable page has no independently authenticated cutoff byte
identity; the source limitation is retained. \citep{exec16_cp_aws}

A generated capacity example makes the distinction operational. Suppose
80\% of customers have baseline recovery probability $0.9$ and a contact
effect $0.02$, while 20\% have baseline probability $0.1$ and effect $0.2$.
At USD~100 per recovery and USD~5 per contact, incremental net values are
$-3$ and $15$. Treating everyone gives only
$0.8(-3)+0.2(15)=0.6$ dollars per customer; contacting the high-increment
group gives $0.2(15)=3$. If capacity permits contacting 10\% of all customers,
randomly selecting half of that group yields $1.5$ dollars per customer.
Ranking by treated response probability instead favors the low-increment
group, whose response probability is $0.92$ rather than $0.3$, and gives
$0.1(-3)=-0.3$ dollars per customer at the same capacity.
\citep{exec16_cp_exact}

The values are stipulated, but the comparison identifies a product requirement:
retain the baseline outcome, incremental effect, cost, assignment probability
and capacity separately. JMP's inspected uplift platform explicitly chooses
partition splits to emphasize treatment differences, using treatment-by-split
interactions for continuous or categorical responses. It is consequently a
relevant targeting comparator. Its older guide's validation and treatment-level
details are scoped to that guide; the current overview cannot establish that
every edition implements the same procedure. Stata's prespecified policy-value
target provides another relevant comparison once the net outcome and population
are aligned. Neither source licenses an automatic inference from a fitted
ranking to an independently validated policy gain. \citep{exec13_ch_jmp_uplift,exec13_ch_stata}

\subsection{Selecting a policy is an additional experiment}

Suppose $m$ candidate policies all have true incremental value zero, and each
evaluation has an independent error equal to $-4$ or $4$ with equal probability.
Every individual estimate is unbiased. Reporting the largest estimate gives
\[
 E\left[\max_{j\leq m}\widehat V_j\right]
 =4\{1-2^{1-m}\}.
\]
The maximum is $-4$ only when all errors are negative and is $4$ otherwise.
Its expectation is zero for one candidate, two for two candidates, and
$127/32$, approximately $3.969$, for eight. A fresh independent mean-zero
evaluation of the selected index has expectation zero. This exact finite
example explains the role of an independent final evaluation; it is not a
proposed universal subtraction rule for a commercial optimizer.
\citep{exec16_cp_exact}

The candidate list, training and tuning records, validation sample and final
policy identity should therefore be part of the comparison. If a product
supports repeated online adaptation, the evaluation must reproduce that
adaptation under an appropriate experiment rather than call it a frozen-policy
test. The best policy among a declared finite list also differs from an
unknown optimum over all feasible strategies. The proposed reporting should
state which value is being estimated and preserve the selection uncertainty
or the independence that justifies its exclusion. \citep{jiang2016,exec16_cp_matrix}

\subsection{Influence diagrams support a different reuse decision}

SMILE's inspected manual represents controllable decisions as nodes with
available states, value nodes with utility tables, and multi-attribute utility
nodes with weights or expressions. It exposes indexing parents, including
earlier decisions and chance observations relevant before a decision. After
inference, its matrices contain policy-indexed probabilities or expected
utilities, depending on node type. GeNIe supplies the associated decision
framing, preference elicitation and value-of-information workflow. These
capabilities are relevant to a bank that can supply a credible finite decision
model. \citep{exec16_cp_bayesfusion}

The proposed reuse choice is a licensed supplied-model comparator. A bridge
must preserve which observations precede which actions, decision-state order,
utility units, the interpretation of indexing parents, and the distinction
between evaluating a supplied policy and optimizing within the model.
Maximizing a supplied utility table cannot validate its transition probabilities
against a selectively logged bank history. Nor does the presence of hidden
chance nodes identify a partially observed causal system from arbitrary data.
The inherited SMILE licensing documentation remains relevant; public manuals
do not grant embedding rights or establish a current price. Porting the vendor
engine is not the proposed route. \citep{exec16_cp_bayesfusion,exec13_smile_license}

\subsection{What the proposed product still needs}

The common local reference revision examined in this probe exposed six
bounded scientific operations; its declared operation schema and estimator
dispatch did not provide a general policy-evaluation service. Three actual local requests naming fixed-history,
learned-policy validation and static-assignment value methods returned
\texttt{unsupported} with no result. This observation is restricted to that
service. The repository also contains separate history-evaluation and
controlled-model research modules; their existence and selected demonstrations
are not equivalent to a qualified common library/CLI/MCP policy interface.
The later fitted servicing and selected-policy studies
(Chapters~\ref{gap29:servicing} and~\ref{gap29:policy-uncertainty}) develop
bounded inferential evidence. The fitted-summary interface demonstration in
Chapter~\ref{gap29:public} is retained separately. None establishes a general
raw-history policy estimator or validates bank logging probabilities.
\citep{exec16_cp_runtime}

The concrete proposal gap is therefore larger than adding another estimator
identifier. A policy request must bind its initial population, physical and
availability clocks, full history variables, action probabilities, observation
and censoring process, discounting, net-reward map, selection/training identity
and final evaluation units. Its response should distinguish a static rule,
a fixed sequence policy, an adapting learner and a supplied-model optimum.
It must retain support failures and uncertainty limits. The documented
commercial workflow features---policy comparison, model export, diagnostics
and decision-model inspection---are useful requirements for this design, while
the scientific meaning is checked independently. \citep{exec16_cp_matrix}

The bounded portfolio covers eleven named families. Stata and JMP support
documented static alternatives; Azure supplies a retired workflow reference;
Amazon supplies a recommendation-evaluation comparator; and BayesFusion
supplies model-based decision analysis. MATLAB's selected official routes were
inaccessible. SAS reinforcement-learning routes and HUGIN's selected routes
did not yield technical originals in that acquisition. The later HUGIN GUI
9.8 manual adds a finite supplied-model decision comparator in
Section~\ref{gap29:software-followup}; its discrete chance-node restriction
and prescribed information sets must remain explicit. It does not supply
off-policy identification from logged histories or qualify a reinforcement
learner \citep{gap29hugin}. The inherited SAS survival, IBM regression,
Mplus dynamic-model and Bayesia causal-model sources cannot settle their
policy-specific capabilities. These entries remain unresolved with concrete
manual/version acquisition routes. No absence of a policy capability is
inferred from an unavailable page. Exact editions, entitlements, prices and
runtime behavior remain open across the prospective commercial trials.
\citep{exec16_cp_matrix}

An economical prototype can begin with the exact mechanisms above and the
existing open-source reference comparisons. Public data can supply realistic
covariates, event clocks and missingness patterns, while generated assignments
and outcomes provide declared causal truth. Compare one static policy,
one history policy and one supplied-model decision on matched eligible
questions; retain a predictive recommendation baseline as an explicitly
different alternative. Freeze final trajectories, candidate policies, metrics,
cost accounting and stopping rules before evaluation. Native commercial
execution follows only for a selected edition with available rights. Bank
experiments become relevant when the prototype identifies a consequential
logging, support, hidden-state or transport uncertainty that public/generated
evidence cannot settle. This route makes the proposal executable while leaving
the unresolved general research questions visible. \citep{exec16_cp_matrix}
\label{exec16:commercial-policy-end}


% SOURCE docs/release/2026-10-29/source/docs/writing-repair/2026-10-29/src/ac94403507ad3147.tex CHAR 0
\section{SMILE and GeNIe: a commercial engine is not a causal evidence contract}
\label{sec:exec24-smile-reuse}

A banking team can use an interactive graphical-model environment to discuss
a lending or servicing decision, and later wish to run the agreed model in
a service. BayesFusion's GeNIe and SMILE address related parts of that
workflow. They should be compared as a connected offering without being
counted as two independent inference engines. The proposal makes two reuse
choices: one for the compiled SMILE engine and its wrappers, and one for the
GeNIe-centered modeling environment and its route to that engine. Both are
reference decisions with conditional integration routes, not acquired
entitlements or verified deployments \cite{exec13_smile_intro,
exec13_smile_license,exec24_smile_choice}.

\subsection{What the retained official documentation establishes}
The retained SMILE introduction declares manual version 2.2.4.R1, built on
27 April 2024; the retained GeNIe identity page declares 5.0.R2, built on
23 November 2024. The primary pages were acquired later and are mutable.
Those declarations identify the inspected manual families; they do not
prove that every byte was served at the 7 October 2026 research cutoff,
or that these are the current installed versions. The previous official
source audits, local bytes and acquisition records remain preserved
\cite{exec13_smile_intro,exec13_genie_readme,exec24_smile_choice}.

The SMILE introduction describes a proprietary C++ library distributed in
compiled form, with Java, Python, R, .NET and COM wrappers. Its C++ manual
expressly leaves most wrapper interoperability to separate documentation.
It describes SMILE as GeNIe's API and relates elementary GUI operations to
engine calls. That relationship supports a shared vendor comparison. It
does not prove exact wrapper-version parity, lossless GUI-to-service export,
or identical interruption, error and state-lifetime behavior across languages
\cite{exec13_smile_intro}.

The inspected licensing page distinguishes academic and commercial
development licenses from server and end-user deployment licenses. It
describes personalized license-key headers and points to the vendor for
terms and pricing. Consequently, an academic evaluation, desktop modeling
session or development license is not evidence of entitlement for a banking
server or redistributed application. The proposal accepts no agreement,
uses no license key and assumes no price. Rights and supported artifact
profiles are inputs to a later controlled integration decision
\cite{exec13_smile_license,exec24_smile_choice}.

\subsection{Discovery output must retain unresolved orientation}
The official SMILE PC interface learns into a \texttt{DSL\_pattern}, with
one node per dataset variable. Its documented controls include a maximum
adjacency, a significance threshold and background knowledge. A successful
call establishes completion under that documented interface; it does not
establish the truth of causal assumptions, the sufficiency of the sample,
or identification of a desired business intervention
\cite{exec13_smile_pc,exec24_smile_choice}.

The inspected pattern reference lists no-edge, undirected and directed edge
types, a conversion to a DAG, and conversion to a network using dataset
variables and outcomes. That conversion assumes matching dataset and
pattern order. The GeNIe pattern editor separately documents graphical
marks and a Create DAG action that randomly chooses directions to produce
an acyclic graph. The GUI description and the narrower inspected API edge
enumeration must not be equated automatically with a standard CPDAG or PAG
representation. Their exact endpoint mapping remains a public conformance
question; an apparent double arrow in a picture is not enough evidence
for a particular latent-confounding interpretation
\cite{exec13_smile_pattern,exec13_genie_pattern,exec24_smile_choice}.

Consider an exact authored binary distribution with
\[
 p(A=0,Y=0)=p(A=1,Y=1)=\frac25,\qquad
 p(A=0,Y=1)=p(A=1,Y=0)=\frac1{10}.
\]
One structural model has a fair action $A$ causing $Y$, with success
probabilities $1/5$ and $4/5$ under actions zero and one. Its intervention
effect is $3/5$. Another has fair $Y$ causing the observed assignment $A$,
with the same conditional agreement probability $4/5$. It generates exactly
the same observed table, but changing $A$ has zero effect on $Y$.
These are two compatible causal explanations, not a claim that they exhaust
the unrestricted identified set. Choosing a DAG orientation selects one
explanation; it does not make the observational evidence distinguish them
\cite{exec24_smile_choice}.

For a hypothetical 2,000-account campaign worth 100 USD per incremental
success and costing 40,000 USD, these two effects imply incremental net
values of 80,000 and minus 40,000 USD. The 120,000 USD difference arises
before any numerical solver error. An agent drafting a diagram must retain
the orientation ambiguity, ask what substantive or experimental information
could resolve it, and record any imposed direction as an assumption.
Generating a visually complete network is not causal validation
\cite{exec24_smile_choice}.

\subsection{Structure conversion and parameter estimation are distinct}
The PC reference says that conversion of its output to a network produces
uniform probability distributions. The GeNIe pattern editor likewise
distinguishes learning parameters from creating only a structure, whose
tables are initially uniform. For a binary outcome, a row $(1/2,1/2)$ can
therefore be an initialization convention rather than an estimated bank
response. A wrapper must carry a parameter-status field and the data or
elicitation provenance for every table. It must refuse a decision-value
claim when a structural scaffold is presented as a fitted model
\cite{exec13_smile_pc,exec13_genie_pattern,exec24_smile_choice}.

This distinction also informs integration with language-model agents.
An agent can propose node labels, information times, candidate edges and
elicitation questions, and can help translate a business objective into a
query. The proposed system must keep those suggestions separate from
estimated tables and licensed engine execution. An attractive diagram,
syntactically valid exported network or uniform default table should not
be promoted by an automated workflow into a supported causal answer.
Human review of substantive assumptions and an explicit scientific query
contract remain necessary parts of the proposed workflow
\cite{exec24_smile_choice}.

\subsection{Temporal models require an explicit information clock}
The retained SMILE dynamic-model tutorial distinguishes initial and temporal
conditional tables, lag order and slice count. It sets temporal evidence
through integer outcome indices and reads beliefs in slice-major order:
the index for slice $t$ and outcome $k$ is $tK+k$, where $K$ is the number
of outcomes. These details are scientific interfaces because the same
numeric array can express different events under another mapping
\cite{exec13_ct_smile_tutorial}.

For an authored three-slice binary example, let the flattened array be
\[
 (0.8,0.2,0.6,0.4,0.3,0.7).
\]
Slice-major decoding gives first-outcome
probabilities $(0.8,0.6,0.3)$. Treating the first three entries as one
outcome's trajectory instead gives $(0.8,0.2,0.6)$. At the second slice,
1,000 hypothetical accounts worth 100 USD per success differ by
40,000 USD in expected value. Tests must therefore vary outcome order,
slice count, lag, variable identity and evidence time rather than merely
checking that an array has the expected length \cite{exec24_smile_choice}.

The retained GeNIe DBN-learning page expressly describes parameter learning
for an already fixed structure and says that its structure-learning
algorithms do not learn dynamic-model structure. It requires time-indexed
column naming and variable/state matching; existing network evidence is
ignored in the described parameter-learning process. That statement is
bounded to this inspected manual page. It supports comparison of the
documented workflow, not an exhaustive claim about every vendor product or
future release. A proposed longitudinal causal product still needs an
assignment model, censoring and treatment-history rules, support,
identification, sensitivity and uncertainty; a fitted DBN alone supplies
none of those by declaration
\cite{exec13_ct_genie_learning,exec24_smile_choice}.

\subsection{Expected utility depends on what was known before the decision}
The supplied-model decision-analysis pages document chance, decision,
value and multi-attribute utility nodes. The SMILE matrices have different
interpretations by node type and are indexed by specified parents:
probabilities for chance nodes, expected utilities for alternatives at
decision nodes, and expected utility for value nodes. Treating every matrix
entry as a probability would corrupt both units and meaning. The adapter
must preserve the decision's information set and the utility's monetary
or other declared scale \cite{exec16_cp_bayesfusion}.

An elementary authored example illustrates the information restriction.
Let a fair binary state determine which of two actions yields 100 USD,
with the other yielding zero. Without observing the state, the best
expected payoff is 50 USD. If a perfect observation is genuinely available
before the action, it is 100 USD and the value of that information is
50 USD. Using an observation recorded after the action as though it were
known before it creates an apparent improvement that was not feasible.
This is why a business-to-diagram assistant must record availability time,
not just a column name or predictive relationship
\cite{exec24_smile_choice}.

The documented utility machinery offers a useful comparator for decision
framing and calculation conditional on a supplied model. It does not, by
itself, establish off-policy identification from logged bank assignments or
sampling-valid uncertainty. The proposed product's additional obligations
are to explain how the model and query are justified, when they are not,
and what public or later bank evidence would change that judgment
\cite{exec16_cp_bayesfusion,exec24_smile_choice}.

\subsection{Two reuse choices and a common public test program}
For the engine decision, reference use is retained and an entitled isolated
worker is conditional on exact C++/wrapper versions, platform support,
deployment rights and scientific fixtures. Direct adoption as a general
causal platform is deferred. Porting the proprietary implementation is not
an available source-reuse route without separate authorization; independently
implementing public methods would be a different project with its own
evidence obligations. Unsupported causal claims and collapsed error types
are rejected, while bounded supplied-model uses remain candidates
\cite{exec13_smile_intro,exec13_smile_license,exec24_smile_choice}.

For the GeNIe/SMILE offering decision, the GUI remains a reference for expert
elicitation, model discussion and decision visualization. Conditional use
requires a documented export/API mapping that preserves unresolved edges,
parameter provenance, temporal semantics and utility units. It must not
silently convert GUI choices into approved server actions. The engine and
GUI choices share a vendor and much evidence; keeping two explicit
decisions reflects their different interfaces and rights, not two qualified
inference capabilities \cite{exec24_smile_choice}.

Public fixtures can test the two equivalent observed models, unresolved
edge export, uniform initialization, shuffled dataset columns, temporal
outcome permutations, decision information timing and typed interruption.
The PC documentation says that a progress callback can interrupt learning
and that elapsed-time checks occur after an iteration. A hard service
deadline therefore needs a separately tested process-control contract,
including safe cancellation and recovery; the documented option alone is
not proof of a strict wall-clock bound
\cite{exec13_smile_pc,exec24_smile_choice}.

The proposed CLI/MCP boundary returns graph class and unresolved marks,
parameter status, target and time information, supported numerical result,
typed refusal or interruption, and complete request provenance. It does
not execute business actions merely because an agent requested inference.
Scientific ownership covers assumptions and identification; adapter
ownership covers mappings and recovery; platform and procurement ownership
cover supported artifacts, entitlements and cost. An illustrative
24,000 USD divided by 800 submitted, 600 completed and 400 qualified
requests gives 30, 40 and 60 USD. Actual quotes, wrapper parity, server
profiles, maintenance and native performance remain unresolved. Hidden
bank assignment and arbitrary causal ambiguity require new information or
justified restrictions. Package execution and independent validation remain to be performed \cite{exec24_smile_choice}.


% SOURCE docs/release/2026-10-29/source/docs/writing-repair/2026-10-29/src/94935a897f851a71.tex CHAR 0
\section{Commercial network tools: separate the query, interface and entitlement}
\label{sec:exec25-bn-commercial-reuse}

BayesiaLab and Bayes Server provide two concrete commercial references for
the proposed banking workflow. The relevant decision is not whether a
graphical product can display a causal diagram. It is which causal query a
selected edition and interface can execute under a supplied model, what
evidence supports that model, and which deployment and review costs remain.
The inspected official documentation is retained as acquired in October
2026. Mutable pages and an undated license text are not immutable snapshots
of the 7 October cutoff, a current quotation, or an accepted bank agreement.
The comparison therefore supports reference choices for named offerings and
conditional prototype designs, with no installed vendor capability claimed.
\citep{exec13_bayesia_api,exec13_bayesia_professional,exec14_ci_release_history,exec14_ci_eula}

\subsection{What the documented workflows contribute}

BayesiaLab's official material distinguishes a desktop modeling environment
from development and production API configurations. Its causal walkthroughs
describe graphical adjustment in a supplied DAG, graph surgery that removes
incoming arcs under an intervention, and likelihood matching that balances
the confounder distribution in the illustrated model. These are useful
comparators for explaining a banking analysis: an analyst can connect the
assumed confounding paths to a concrete intervention calculation instead of
presenting an unexplained effect number. They do not establish that an
automatically learned banking graph has the right causal directions or that
its unobserved causes have been resolved. The selected examples also do not
provide an exhaustive inventory of supported identification algorithms.
\citep{exec13_ci_bayesia_criteria,exec13_ci_bayesia_surgery,exec13_ci_bayesia_matching}

The Bayesia Engine page describes a REST service for inference with an
existing Bayesian network. It separately lists Java configurations for
modeling, learning and inference. A proposal adapter must bind the actual
configuration and query contract; a REST endpoint's existence does not show
that it exposes desktop learning or every causal operation. The proposed
library currently lacks a qualified export/import and inference route for
these profiles. It should preserve node identities, state order, evidence,
intervention status, parameter provenance and units before comparing outputs.
Where a graphical adjustment explanation is part of the analyst journey,
the report should retain its assumptions rather than return only an API
number.~\citep{exec13_bayesia_api,exec13_ci_bayesia_surgery}

Bayes Server's guide explicitly separates ordinary evidence from
interventions and demonstrates an abduction--intervention workflow. In its
individual example, factual evidence is entered, characteristic variables
are inferred, the original evidence is removed, and the alternative
intervention is applied. The guide distinguishes calculations available
from version 9.4 from helper functions available from version 10. It also
describes several identification approaches and recognizes that a valid
criterion need not exist. The retained release history lists version
11.4.0.0 dated 17 April 2026. This is a fact about that acquired history; it
does not establish a currently installed build or parity across all APIs.
\citep{exec13_ci_bayesserver_causality,exec14_ci_release_history}

These workflows expose two proposal gaps. First, the library has not
demonstrated an analyst experience that traces a supplied model through
adjustment, intervention and an intelligible explanation. Second, it has
not qualified a stateful factual-to-counterfactual workflow against these
commercial interfaces. A generic graph schema and a callable numerical
engine alone do not close either gap. They motivate separate public
prototype tasks, using small stipulated models with exact answers before
an expensive bank study.~\citep{exec25_bn_commercial_choice}

\subsection{Intervention agreement does not settle an individual counterfactual}

An exact authored example clarifies what the interface must preserve. Let
$X$ be randomized offer assignment and let $U$ be an unobserved binary
customer characteristic, independent of $X$, with both Bernoulli$(1/2)$.
Consider two stipulated response laws, with $\oplus$ denoting exclusive-or:
\[
 A:\quad Y=U,\qquad B:\quad Y=X\oplus U.
\]
These are alternative mathematical models, not actual customers or fitted
vendor networks. No faithfulness or monotonicity assumption is imposed.
Enumerating the four equally probable $(X,U)$ states gives
\[
 P_A(X=x,Y=y)=P_B(X=x,Y=y)=\tfrac14
 \quad(x,y\in\{0,1\}).
\]
Under either intervention, the marginal response remains one half in both
models, so both average treatment effects are zero:
\[
 E_A(Y\mid\operatorname{do}(X=x))
 =E_B(Y\mid\operatorname{do}(X=x))=\tfrac12.
\]
Now consider a customer factually observed with $X=0,Y=0$. In either
stipulated model, those facts imply $U=0$. Holding that inferred
characteristic fixed and setting $X=1$ gives
\[
 P_A(Y_1=1\mid X=0,Y=0)=0,\qquad
 P_B(Y_1=1\mid X=0,Y=0)=1.
\]
Thus matching both observational probabilities and all these population
intervention means does not determine this individual counterfactual. The
models couple potential responses differently. Additional structural
restrictions could exclude a model; the numerical interface must record
those restrictions rather than silently manufacture them. This example
establishes ambiguity in a class containing both stipulated models, not a
general sharp identified set or a defect in either vendor.
\citep{exec25_bn_commercial_choice}

For a banking request, a population offer effect, a contribution calculation
and an individual's alternative history are therefore separate result
types. The inspected Bayesia contribution material describes factual and
intervention contrasts under expert confounder assumptions. Its title alone
does not make every such contrast the particular cross-world quantity in
the example. The Bayes Server abduction workflow provides a concrete
reference for retaining characteristic-variable information, but the truth
of its supplied structural coupling remains a modeling obligation. A
proposed CLI or MCP response should name the target, factual evidence,
intervention, retained exogenous information and the assumptions supporting
their combination. When that information is absent, an explicit unresolved
counterfactual is more defensible than an unqualified number.
\citep{exec13_ci_bayesia_counterfactual,exec13_ci_bayesserver_causality,exec25_bn_commercial_choice}

\subsection{Deployment units and commercial evidence}

The Bayesia Engine documentation defines a production instance as a loaded
network model ready for an operation and describes licensing by simultaneous
instances. It distinguishes this from the physical server count. This has
a direct implication for a prototype scheduler: record load and unload
intervals, peak occupancy, waiting time and all model transitions. A
developer seat, a model file, a host and a concurrent loaded instance are
different accounting units. The selected documentation provides a model
for investigating them; actual entitlement and pricing require the
selected configuration and applicable agreement.
\citep{exec13_bayesia_api,exec13_bayesia_professional}

For a simple authored schedule, model $A$ is loaded on $[0,5)$, $B$ on
$[3,7)$ and $C$ on $[5,9)$. With half-open intervals, occupancy is
\[
 n(t)=\sum_{m\in\{A,B,C\}}\mathbf1
       \{\ell_m\le t<u_m\},\qquad \max_t n(t)=2.
\]
There are three distinct models but a peak of two loaded models. The result
is independent of whether those intervals occur on one physical host or
several. It is an accounting illustration, not a determination that a
particular license permits the schedule. A serialized schedule might reduce
concurrency while increasing latency; both costs belong in the comparison.
\citep{exec25_bn_commercial_choice}

Bayes Server's retained purchase page presents version-11 developer and
desktop options, while its linked license text contains executable-use,
transfer and benchmark-disclosure conditions. These documents must be
reconciled with the actual selected order and any addendum before a
deployment or externally reported native comparison is promised. The
purchase description alone does not settle the applicable rights. No
agreement is accepted and no vendor is contacted by this proposal review.
The rights question remains a local prerequisite for that native trial;
it does not prevent independent mathematical examples or continued
proposal research.~\citep{exec14_ci_purchase,exec14_ci_eula}

\subsection{Two reference decisions and their qualification routes}

For BayesiaLab, retain the desktop explanation and specified Engine
configurations as separate reference profiles within one offering.
Conditional integration would first select the exact desktop/API artifact
and licensed deployment, then test model export, state ordering, evidence
versus intervention, adjusted versus unadjusted targets, default versus
estimated parameters and failure recovery. For Bayes Server, retain the
documented intervention and abduction workflows as references. Conditional
integration must additionally preserve characteristic-variable evidence and
the order of factual conditioning, abduction, evidence removal and
intervention. Neither choice authorizes direct adoption or a source port.
An independent implementation of public mathematics is a separate option
with its own evidence and equivalence obligations, not an implicit right
to copy a proprietary implementation.~\citep{exec25_bn_commercial_choice}

Both public qualification plans would compare a hand-enumerated model, the
selected native profile when permitted, and the proposed adapter under the
same target and information. Scientific checks include exact intervention
probabilities, the counterfactual distinction above, unsupported hidden
confounding and correctly refused queries. Runtime checks include
variable/state permutations, graph and parameter round trips, malformed
requests, interrupted sessions, stale models, duplicate requests and
license-server or service unavailability. A successful transport response
cannot substitute for scientific correctness. These are proposed checks;
no vendor binary or service was executed for this review.
\citep{exec25_bn_commercial_choice}

Scientific ownership covers the supplied model, target and uncertainty.
Adapter ownership covers state, graph, evidence and failure conversions.
Platform and procurement ownership covers the supported artifact, network
conditions, deployment rights and pricing. Review ownership covers
independent oracles, permitted reporting and complete denominators. These
are proposed roles requiring actual owners at prototype approval.
Maintenance must follow both numerical behavior and changes to commercial
profiles; a successful historical source review is not continuing support.
\citep{exec25_bn_commercial_choice}

Suppose, purely illustratively, that model preparation costs USD\,2,400,
all requests USD\,1,200, review USD\,3,600 and allocated maintenance
USD\,800. The USD\,8,000 total costs USD\,40 per request across 200
attempts, USD\,50 per completion across 160 completions, and USD\,80
per admissible result across 100 results. These are exact authored
denominators, not product prices or observed bank value. A comparison
should preserve failures, refusals, model-loading costs and human review
rather than reporting only the cost of a completed inference call. The
two reference choices can be completed now; adoption awaits selected
rights, exact profiles and matched scientific and operational evidence.
\citep{exec25_bn_commercial_choice}
\label{sec:exec25-bn-commercial-reuse-end}


% SOURCE docs/release/2026-10-29/source/docs/writing-repair/2026-10-29/src/c81e068ace0033cb.tex CHAR 0
\section{Stata: an integration decision must name the command and estimand}
\label{sec:exec24-stata-reuse}

A universal bank may already use Stata for treatment effects, panel
evaluation, survival analysis and macroeconomic models. That breadth makes
it an important commercial comparator, but the product name does not define
a single scientific contract. The same option label can have different
meanings across commands, and a returned estimate may target a population,
a treated subgroup, a cohort-period cell or a supplied structural model.
The proposed reuse route is therefore a command-specific licensed worker,
with Stata retained now as a reference. A blanket adoption decision would
conceal the very assumptions that the causal proposal is meant to expose
\cite{exec13_ci_stata_causal_intro,exec18_gp_stata_xthdid,
exec22_state_stata_sspace_post,exec24_stata_choice}.

\subsection{The artifact and entitlement are part of the comparison}
The earlier surveys retain official command manuals, version declarations,
local originals and inspected technical anchors across the relevant areas.
Some manuals use release-19 example routes or carry 2025 copyright and
2026 PDF metadata; some features are expressly marked StataNow. These
facts must remain distinct from an exact installed edition and from proof
of the bytes available at the 7 October 2026 research cutoff. The present
decision preserves the earlier dated-source limits and re-inspects selected
target, clustering, pooling and prediction contracts. It does not present
a fresh exhaustive inventory of every Stata extension
\cite{exec20_med_stata_multiple,exec18_gp_stata_post,exec24_stata_choice}.

The retained license-options page distinguishes single-user, concurrent
network and site arrangements, as well as annual and perpetual terms.
A network license's described concurrent users and installation scope do
not by themselves establish an entitlement for an agent-facing service,
unattended shared worker or redistributed application. Exact supported
edition, process topology, concurrency, rights and price remain procurement
and platform inputs. No license agreement is accepted, binary run or
commercial quote invented in this proposal \cite{exec13_ivrd_license,
exec24_stata_choice}.

Bundled commands must also be distinguished from separately maintained
Stata packages. The preserved reviews of \texttt{rdrobust},
\texttt{rddensity} and the historical \texttt{synth} interface carry their
own pinned identities, source and notices. A Stata license does not
establish their scientific equivalence to bundled commands or remove their
dependency and maintenance obligations. Conversely, an open package's
source permission does not supply the commercial runtime entitlement
\cite{exec13_ivrd_rdrobust,exec13_ivrd_rddensity,exec19_synth_stata_source,
exec24_stata_choice}.

\subsection{Treatment effects and heterogeneous targets}
The treatment-effects introduction documents average treatment effects,
effects among treated subjects and potential-outcome means, with regression
adjustment, inverse-probability weighting, augmented methods and matching
families. The interpretation depends on stated sampling, conditional
independence and support assumptions. A fitted propensity score or
successful command does not verify the absence of omitted confounding.
The adapter must preserve the outcome, treatment levels, target population,
covariates and estimation sample as part of the result
\cite{exec13_ci_stata_causal_intro}.

The heterogeneous-effects survey distinguishes conditional average effects,
prespecified groups and groups learned from the data. The documented IATE
label is a covariate-conditional average, not direct observation of an
individual's two potential outcomes. Its nuisance learners, cross-fitting,
honesty and policy-evaluation options have specific contracts. Random
observation folds do not automatically preserve a bank's repeated-customer,
branch or time dependence. A proposed worker must declare the fold and
sampling regime separately from the learner name
\cite{exec13_ch_stata,exec24_stata_choice}.

The same discipline applies to panels. The inspected
\texttt{xthdidregress} manual distinguishes panel data from repeated
cross-sections, defines treatment-assignment groups, and provides
never-treated or not-yet-treated controls. Its \texttt{robust} variance
option is documented as equivalent to clustering on the assignment-group
variable. That is a command-specific statement, not a universal translation
rule for an option called robust elsewhere
\cite{exec18_gp_stata_xthdid}.

\subsection{Aggregation weights can change the investment decision}
Let $\theta_{gt}$ denote a cohort-period effect and let $w_{gt}$ be its
declared aggregation weight. The target is
\[
 \theta(w)=\sum_{g,t} w_{gt}\theta_{gt},\qquad
 w_{gt}\geq0,\quad\sum_{g,t}w_{gt}=1.
\]
The retained postestimation manual offers overall, dynamic, time and cohort
aggregation. Its default time-cohort weights use the subset for each
cohort-period effect, whereas cohort weights pool subsets across periods.
The target and the contributing samples must therefore travel with an
aggregated number. The manual also distinguishes pointwise intervals from
an optional simultaneous construction, with estimator-specific restrictions
\cite{exec18_gp_stata_post}.

An exact authored example has two effects, 10 and minus 5 USD per account.
Weights $(4/5,1/5)$ give 7 USD, while weights $(1/5,4/5)$ give minus
2 USD. For a hypothetical 10,000-account target these are 70,000 and minus
20,000 USD before implementation costs. Neither arithmetic result is a
software defect: they answer different weighted-population questions.
Dropping weights during export would nevertheless turn a correct estimate
into the wrong investment input. This example does not claim that either
weight vector reproduces a particular native default
\cite{exec24_stata_choice}.

Dependence changes uncertainty independently of those target weights.
For an authored illustration, suppose there are ten independent branches,
each with one mean-zero, unit-variance shock repeated identically for its
100 accounts. The mean of all 1,000 account records is the mean of ten
shocks, with variance $1/10$. Treating the records as independent would
give $1/1000$, understating the variance by 100 and the standard error by
10. This is an exact known-covariance example, not evidence that a
ten-cluster asymptotic interval is well calibrated. It explains why
assignment level, unit identity and the proposed finite-sample diagnostic
must survive any CLI or MCP conversion \cite{exec24_stata_choice}.

\subsection{Longitudinal breadth does not remove regime restrictions}
The selected survival-treatment documentation provides targets such as
potential mean survival times and ATE/ATET under its stated models. It
expressly excludes delayed entry, time-varying covariates and multiple-record
data in the inspected family. Its censoring and overlap conditions are
substantive restrictions. This creates a concrete comparison with a bank's
repeated servicing and changing-treatment histories: such histories need
a separately qualified longitudinal method, not merely a survival output
renamed as sequential causal inference
\cite{exec14_survival_stata,exec24_stata_choice}.

The dynamic-panel manuals likewise define a particular regime. The
Arellano--Bond command targets many panels and few periods with specified
moment and innovation assumptions, and exposes lag-instrument and
one-/two-step variance choices. A command's serial-correlation diagnostic
must be interpreted on the differenced errors it actually tests. If level
innovations $\epsilon_t$ are independent with variance one, then
\[
 \operatorname{Var}(\epsilon_t-\epsilon_{t-1})=2,\qquad
 \operatorname{Cov}(\epsilon_t-\epsilon_{t-1},
                    \epsilon_{t-1}-\epsilon_{t-2})=-1.
\]
The resulting first-lag correlation is minus one half even under those
independent level innovations; the second-lag covariance is zero. A generic
rule rejecting any detected first-order correlation would misread this
design. Conversely, non-rejection is not proof of every instrument moment
\cite{exec19_dynamic_stata_ab,exec24_stata_choice}.

Related manuals make further distinctions that the product must preserve:
Granger restrictions concern predictive lag content; structural VARs need
declared restrictions or valid external instruments; system-GMM adds its
own moment assumptions. The proposed library cannot infer causal shock
identity, exclusion or validity from the fact that a numerical estimate
converged \cite{exec13_ct_stata_granger,exec13_ct_stata_svar,
exec13_ct_stata_ivsvar,exec19_dynamic_stata_system,exec24_stata_choice}.

\subsection{Missing-data pooling is part of the uncertainty contract}
The multiple-imputation documentation separates completed-data estimation
from pooling. For $M$ imputations, estimates $\widehat Q_m$, within-imputation
variances $U_m$ and mean $\overline Q$, the familiar scalar components are
\[
 W=M^{-1}\sum_m U_m,\qquad
 B=(M-1)^{-1}\sum_m(\widehat Q_m-\overline Q)^2,\qquad
 T=W+(1+M^{-1})B.
\]
Valid use requires appropriate imputation and completed-data inference;
the retained substantive discussion does not allow observed data alone to
resolve unrestricted missing-not-at-random behavior
\cite{exec21_transport_stata_estimate,exec21_transport_stata_substantive}.

With two authored estimates zero and two, each having within variance one,
the pooled estimate is one, $W=1$, $B=2$ and $T=4$. Reporting only the
within-imputation variance would reduce the standard error from two to
one. This arithmetic witness explains the required fields; it does not
justify an actual imputation model. The manual's default stops on failed
individual estimation or varying estimation samples. Explicit override
options change those behaviors. A worker must preserve every failure,
sample difference and override instead of silently pooling whichever
results happened to succeed \cite{exec21_transport_stata_estimate,
exec24_stata_choice}.

The nonlinear scale of a reported target matters as well. Pooling
coefficients and then exponentiating is not generally equivalent to
averaging already exponentiated estimates. Command, estimand, transformation
and covariance scale must remain explicit in the response. Chained
imputation's conditional models and iteration settings also require their
own diagnostic and compatibility argument, rather than being accepted
because they produced a complete table
\cite{exec21_transport_stata_estimate,exec21_transport_stata_chained,
exec24_stata_choice}.

\subsection{Prediction time and posterior error are different quantities}
The state-space postestimation manual distinguishes one-step prediction
using past information, filtering using past and contemporaneous data, and
smoothing using all sample information. It also states that the reported
prediction RMSE reflects disturbance variability but excludes parameter
estimation error. These distinctions must remain visible when a bank
compares retrospective risk reconstruction with a decision made in real
time \cite{exec22_state_stata_sspace_post}.

In an authored Gaussian example, a static state $S$ has prior $N(0,1)$.
A later measurement $Y=S+\eta$, with independent $\eta\sim N(0,1)$, is
observed as two. Before that observation the conditional mean is zero;
after it the posterior mean is one and variance one half. A retrospective
smoothed estimate can be appropriate for reconstruction, but it cannot
be substituted into the earlier decision's information set. This is a
finite information-clock example, not an installed Stata result
\cite{exec24_stata_choice}.

The Bayesian manual separately documents likelihoods, priors, adaptation
and diagnostics, and cautions that permitted model combinations need not
produce a proper posterior. Monte Carlo error is distinct from posterior
uncertainty, which is itself distinct from repeated-sampling causal
uncertainty. Saved frequency-weighted chain representations and parameter
names must survive export. An agent can help construct and explain a model,
but cannot turn mixing diagnostics into a proof of causal identification
\cite{exec22_state_stata_bayesmh,exec24_stata_choice}.

\subsection{A broad comparison with narrow executable contracts}
The previous source surveys cover additional relevant routes without
granting them common assumptions. Mediation distinguishes one- and
two-mediator commands, variable-type restrictions and causal ignorability.
Spatial regression depends on the weights matrix and its normalization;
matching a graph-shaped input does not establish the bank's interference
mechanism. Structural-equation fitting is evidence for a supplied-model
path, not automatic graph discovery. These command-specific comparisons
remain the basis for deciding which bounded worker would add value
\cite{exec13_ci_stata_mediate,exec20_med_stata_multiple,
exec20_network_stata_spregress,exec13_stata_sem,exec24_stata_choice}.

The reference decision retains these documented capabilities and their
limitations. Conditional integration would run an explicitly entitled
command and artifact in an isolated worker, returning target, sample,
weights, clusters, model assumptions, uncertainty components, diagnostics
and typed failures. Direct adoption as a universal causal platform is
deferred. Porting proprietary implementation is not an available reuse route
without separate authorization; public method reimplementation and
separately licensed extensions have different review obligations.
Unsupported regime or guarantee claims are rejected rather than patched
by relabeling the output \cite{exec13_ivrd_license,exec24_stata_choice}.

Public tests can begin with finite targets and known dependence, staggered
treatment examples, missing-data pooling, information-time fixtures and
failure-preserving batch execution. Scientific ownership fixes the estimand
and assumptions; statistical ownership fixes sampling and covariance;
adapter ownership fixes command versions, sample/weight maps and failure
recovery; platform and procurement ownership fix support and rights.
Hypothetical expenditure of 30,000 USD over 1,000 submitted, 750 completed
and 500 qualified requests gives 30, 40 and 60 USD. Actual prices, runtime,
supported service profiles and institutional value remain unmeasured.
Hidden assignment, unsupported counterfactual responses and arbitrary
dependent-data uncertainty require information or justified restrictions.
Package execution and independent validation remain to be performed \cite{exec24_stata_choice}.


% SOURCE docs/release/2026-10-29/source/docs/writing-repair/2026-10-29/src/e7bf0dee3e069468.tex CHAR 0
\section{JMP: treatment targets, selected subgroups and competing events}
\label{sec:exec24-jmp-reuse}

A bank analyst considering a retention campaign may want an overall
treatment effect, a subgroup worth targeting, or a prediction of default
before another event such as prepayment. JMP's retained documentation
provides distinct relevant workflows. Comparing them requires preserving
the estimand and edition rather than treating a causal-treatment report,
an uplift tree and a survival display as interchangeable evidence. The
proposal retains JMP as a commercial reference with conditional analyst
and worker routes; no native capability or entitlement is established
\cite{exec13_ch_jmp_causal,exec13_ch_jmp_uplift,exec14_survival_jmp,
exec24_jmp_choice}.

\subsection{Declared editions and the inspected technical boundary}
The causal-treatment pages are retained from the JMP Pro 19.2 help family
and display 13 January 2026 publication dates. The uplift overview belongs
to 19.0 and is dated 15 July 2025; the more detailed retained Consumer
Research guide is version 16. The supplied-model instrumental-variable
discussion is from the September 2025 version-19 Multivariate Methods
guide, while the selected survival pages are 18.1 and dated 27 June 2024.
Later acquisition of those mutable pages is not proof of exact bytes at
the 7 October 2026 cutoff. Older-guide behavior is not automatically
verified for a newer installed product
\cite{exec13_ch_jmp_causal,exec13_ch_jmp_uplift,exec13_ivrd_jmp,
exec14_survival_jmp}.

The causal help contains equations rendered as images. The retained
inspection uses the surrounding substantive prose, not an invented
transcription of those equations. The finite equations and examples below
are independently authored and identified as such. No additional image
review is claimed. Exact installed edition, supported automation profile,
commercial rights and prices remain separate questions; a help page is
neither a license nor evidence of an executed service
\cite{exec24_jmp_choice}.

\subsection{An average causal effect has a specified target population}
The documented Causal Treatment personality is available in JMP Pro and
supports binary or continuous treatment and outcome. For binary treatment,
the prose distinguishes average treatment effect from average treatment
effect among the treated. It documents IPWR, regression adjustment and
augmented inverse-probability methods using specified logistic propensity
and linear or logistic response models. For continuous treatment, the
reported targets are mean potential outcomes. These are useful concrete
capabilities to compare, but their interpretation depends on positivity,
consistency and conditional exchangeability
\cite{exec13_ch_jmp_causal}.

The assumption page explicitly distinguishes assumptions about the
data-generating process from empirical checks on a sample. A propensity
plot can reveal observed overlap problems; it cannot establish that every
unmeasured cause was included or that a campaign code corresponds to the
same intervention for every account. A business-to-causal assistant should
therefore record treatment versions, assignment information and alternative
causal explanations alongside any generated model specification. A fitted
response surface is not a substitute for that argument
\cite{exec13_ch_jmp_causal,exec24_jmp_choice}.

Targeting also requires preserving treatment-level coding. The version-16
uplift guide says that, with more than two treatment levels, the first
level is compared with a combination of all remaining levels. It also
documents value-order conventions and validation-column roles. These are
important historical contracts to test in the intended version, not
defaults that can be silently carried forward
\cite{exec13_ch_jmp_uplift}.

An authored randomized example has success probabilities $1/5$, $3/5$
and $9/10$ for a baseline and two distinct offers. If the pooled alternative
uses the latter offers with weights $(3/4,1/4)$, its contrast with baseline
is $19/40$. Reversing the mixture weights gives $5/8$. For 1,000
hypothetical accounts worth 100 USD per incremental success, the values
are 47,500 and 62,500 USD before offer costs. Both are valid mixture
contrasts under the stated experiment, but neither is automatically the
effect of a particular offer. The adapter must retain the first level,
pooled levels, their weights and the business action actually proposed
\cite{exec24_jmp_choice}.

\subsection{A selected uplift split needs its own evaluation law}
The retained uplift overview describes partition splits based on
treatment-by-split interactions, using linear models for continuous
responses and logistic models for categorical responses. It documents a
Monte Carlo adjustment intended to address predictor-selection bias from
different numbers of levels or candidate models. The older guide also
describes validation-based split selection. These features should be
acknowledged as documented controls, without inferring a general
post-selection interval or dependent-data guarantee
\cite{exec13_ch_jmp_uplift}.

A simple authored selection example explains the unresolved statistical
question. Suppose four independent candidate scores each equal minus one
or plus one with equal probability and have true mean zero. The selected
maximum is minus one only when all four scores are negative, which has
probability $1/16$. Its expected value is therefore
\[
 E[\max_j S_j]=\frac{15}{16}-\frac1{16}=\frac78.
\]
Selecting the largest score creates optimism even though no candidate
has a positive mean. This calculation does not test JMP's adjustment or
claim that it produces this bias. It specifies why the proposed comparison
needs a held-out or otherwise justified evaluation law for the complete
selection process, including predictor search, tree depth and policy choice
\cite{exec24_jmp_choice}.

In a bank, splits may be selected repeatedly across branches, customer
histories or calendar periods. Public fixtures should separate training,
selection and evaluation units according to the declared dependence regime.
They should retain all attempted splits and policies, treatment allocation,
costs and failure counts. A language-model agent may explain candidate
segments, but should not convert a selected in-sample difference into a
qualified intervention recommendation without the corresponding evidence
\cite{exec24_jmp_choice}.

\subsection{A continuous density model does not create empirical support}
The continuous-treatment prose specifies a normal conditional treatment
density with parameters estimated through linear regression. It describes
a stabilized-weight marginal structural model and a generalized propensity
score response model. These are particular modeling choices, not evidence
of identification for every dose or treatment distribution. A density
model assigning a positive number outside the observed range can conceal
an unsupported policy extrapolation if the wrapper records only that number
\cite{exec13_ch_jmp_causal,exec24_jmp_choice}.

For an authored bounded outcome, suppose treatments are observed only on
$[0,1]$. Two response laws agree there at $\mu(a)=a/2$. Beyond one, let
the first continue to $\mu_1(2)=1$, while the second declines according to
$\mu_2(a)=1-a/2$ and reaches $\mu_2(2)=0$. Both laws fit the same
observed support but imply opposite responses at the proposed dose two.
A parametric continuation chooses a restriction; it does not reveal the
missing response from the observed sample. Public tests can check that a
worker reports the support boundary and the chosen extrapolation model.
Actual off-support bank response requires new information or a justified
restriction \cite{exec24_jmp_choice}.

\subsection{Competing-event displays are not automatically removal interventions}
The inspected competing-causes page describes estimating each cause's
survival distribution by treating other events as censored, under a stated
independent-causes premise. It offers an Omit Causes display that
recalculates survival estimates to illustrate a situation in which selected
causes are no longer hazardous. Those documented operations do not establish
what would happen if a bank changed a prepayment policy or prevented a
particular competing event
\cite{exec14_survival_jmp}.

An authored finite example clarifies the distinction. Half of accounts
have a competing event at time one and latent default at time two; the
other half default at time two before their competing event at time three.
The observed fraction defaulting before a competing event is one half.
Under a specifically defined intervention that removes the competitor
while leaving each latent default time unchanged, all accounts default
by time two. Other intervention mechanisms could also change default
times. The observed competing-event record alone does not select the
unchanged-latent-time intervention. This example is not a native JMP
calculation and does not assume that its display uses this discrete model
\cite{exec24_jmp_choice}.

The proposed product therefore needs distinct labels for observed
cumulative incidence, cause-specific or net survival, and an explicitly
defined counterfactual removal target. Event coding, censoring, risk sets,
time units and dependence assumptions must be preserved. A plotting or
simulation option is a useful analytical interface; it does not supply the
causal intervention mechanism by itself
\cite{exec14_survival_jmp,exec24_jmp_choice}.

\subsection{The supplied-model IV route and remaining scope}
The retained version-19 multivariate guide documents model-implied
instrumental-variable estimation with equation-level diagnostics and
instrument lists. Its instruments follow a supplied structural model;
those declarations do not make invalid exclusion assumptions true or
establish unrestricted graph discovery. This is another conditional
comparison route, separate from causal-treatment regression and uplift
partitioning \cite{exec13_ivrd_jmp}.

The present evidence does not establish complete panel, sequential-treatment,
latent-discovery or agent-service contracts for every JMP product and
extension. That is a coverage limit, not a claim that the vendor lacks all
such features. The next survey should identify the intended edition,
official API or automation interface, supported artifact and entitlement,
then acquire the concrete missing technical pages. Image-only equations
need a fresh permitted visual inspection before claims relying on them
are admitted \cite{exec24_jmp_choice}.

Reference use is retained for the documented analytical workflows.
Conditional integration may use an entitled analyst workflow or isolated
worker after public target, coding, support, selection and event-time
fixtures pass. Direct adoption as the full causal product is deferred.
Porting a proprietary implementation is not an authorized open-source
reuse route; independent implementation of public methods is a separate
proposal. Unsupported causal, off-support or post-selection guarantees
are rejected without rejecting every bounded use of the software
\cite{exec24_jmp_choice}.

Scientific ownership covers intervention, target, assumptions and loss;
statistical ownership covers selection and dependence; adapter ownership
covers treatment order, mixture weights, events and typed refusals; platform
ownership covers supported automation, rights and complete costs. A
hypothetical 21,000 USD over 700 submitted, 525 completed and 350
qualified requests gives 30, 40 and 60 USD respectively. Actual prices,
support and bank value remain unmeasured. Package execution and independent validation remain to be performed
\cite{exec24_jmp_choice}.


% SOURCE docs/release/2026-10-29/source/docs/writing-repair/2026-10-29/src/7674c5358f19b3c3.tex CHAR 0
\section{SAS: preserve the boundary between statistical procedures and services}
\label{sec:exec24-sas-reuse}

A universal bank's SAS estate may contain statistical procedures, historical
econometric models, a Viya platform and open clients that call particular
services. These are related components with different scientific and operating
contracts. The proposal retains the documented workflows as references and
defines conditional integration at the procedure and service level. A Python
client's availability cannot establish that a bank owns the required server
component, that a supplied structural model is identified, or that an
asynchronous job returned a usable causal result
\cite{exec13_latent_sas_calis,exec23_workflow_sasctl,
exec23_workflow_delegates_sasctl,exec24_sas_choice}.

\subsection{Edition and component identity constrain the comparison}
The retained scientific evidence includes explicitly archived SAS/STAT 14.1
and SAS/ETS 14.1 pages. The latter are adjacent econometric comparisons,
not a declaration that every SAS/STAT or Viya entitlement includes them.
The selected spatial-regression guide declares 2023.10 and carries an
August 2024 document date. These identities remain separate from a current
installation, supported server release and exact original bytes at the
7 October 2026 cutoff. Selected HTML equations are image-based; inspected
prose supports the vendor claims, while the finite equations below are
independently authored
\cite{exec13_ct_sas_causal,exec13_ivrd_sas,
exec20_network_sas_latest_spatial,exec24_sas_choice}.

The open sasctl client is inspected at the retained 26 August 2026 revision.
Its README documents Python-module and command-line use. Its own root notice
declares Apache License 2.0; this is not an entitlement to SAS/STAT, SAS Viya
or every server API. The separate REST sample repository is explicitly
described as a preview. Its Job Execution example's document version and
historical update date are not a supported server compatibility matrix
\cite{exec23_workflow_sasctl,exec24_workflow_contracts_sas}.

\subsection{A fitted structural model preserves its supplied assumptions}
The CALIS overview documents observed and latent structural models,
measurement error, path specifications, constraints and reciprocal relations.
It distinguishes normal-theory maximum likelihood, Satorra--Bentler adjusted
inference, and an outlier-downweighting option that the procedure calls
robust. Those alternatives should not be collapsed into one flag. The
documented full-information missing-data route assumes missing at random;
a successful optimizer does not verify that assumption
\cite{exec13_latent_sas_calis}.

An authored Gaussian example shows why covariance fit does not settle an
intervention direction. For independent standard-normal disturbances, one
model is
\[
 X=U,\qquad Y=\tfrac12X+\sqrt{3/4}\,V.
\]
A second model reverses the direction:
\[
 Y=U',\qquad X=\tfrac12Y+\sqrt{3/4}\,V'.
\]
Both have zero means, unit variances and covariance one half. Setting $X=x$
in the first model changes the mean of $Y$ at rate one half; setting it in
the second changes that mean at rate zero. This is an exact model witness,
not an executed CALIS comparison. A generated path diagram must retain the
business mechanism and identifying restrictions that distinguish the two
models, rather than accepting whichever orientation gives a convenient
story \cite{exec24_sas_choice}.

The SYSLIN documentation similarly treats instruments as specified inputs:
the listed variables enter first-stage regressions for the relevant
estimators. Declaring those variables does not prove exclusion or their
relationship to omitted disturbances. The VARMAX CAUSAL statement tests
specified Granger restrictions; its label does not automatically identify
a bank intervention. The dynamic simultaneous-equations representation is
another declared model representation, not evidence that an unrestricted
shock interpretation has been learned
\cite{exec13_ivrd_sas,exec13_ct_sas_causal,exec13_ct_sas_structural}.

\subsection{Event histories require the information available at each time}
The PHREG overview provides a concrete strength relative to a generic
single-record treatment interface: time-dependent explanatory variables,
counting-process input, recurrent events and left truncation are explicitly
discussed. It also discusses clustered failure times and different methods
for tied times. The adapter should preserve entry and exit times, event
types, subject and cluster identity, time units and the selected method
\cite{exec14_survival_sas}.

Support for a changing covariate does not establish that the covariate was
known when the bank acted, or that conditioning on it identifies a dynamic
treatment strategy. In an authored four-account example, let the eventual
event indicators be zero, zero, one and one. Copying those future outcomes
into an alleged baseline covariate gives perfect prediction by construction.
It provides no feasible baseline decision rule. A public fixture should
reject this timestamp violation even when its numerical fit is excellent.
With treatment-dependent intermediate variables, the proposed scientific
contract must additionally specify the longitudinal intervention and the
assumptions supporting its estimator
\cite{exec24_sas_choice}.

The selected spatial-regression documentation adds a further target
distinction. Dependence through a supplied spatial matrix, direct and
indirect effects, and the banking action that changes a network exposure
need separate definitions. The existing spatial comparison retains those
conditions. No product-wide permission to infer causality follows from
having both a survival procedure and a spatial procedure available
\cite{exec20_network_sas_latest_spatial,exec24_sas_choice}.

\subsection{Missing outcomes remain a scientific sensitivity question}
The MI and MIANALYZE manuals distinguish imputation, completed-data analysis
and pooling of estimates with their uncertainty. The assumptions pages
state missing at random and distinguish the data-model parameters from
the missingness-model parameters. The sensitivity discussion treats
pattern-mixture and selection-model factorizations separately; the
inspected edition implements the described pattern-mixture route and
explicitly says its MI procedure does not provide the selection-model
approach discussed there. These are scoped edition-specific statements
\cite{exec21_transport_sas_mi_intro,exec21_transport_sas_mianalyze_intro,
exec21_transport_sas_assumptions,exec21_transport_sas_sensitivity}.

For an authored bounded outcome, suppose exactly half the population is
observed, its observed mean is $1/5$, and every outcome lies between zero
and one. If the missing half has mean $m$, the population mean is
\[
 \mu=\tfrac12\cdot\tfrac15+\tfrac12m,
 \qquad 0\leq m\leq1.
\]
The observed information permits means from $1/10$ to $3/5$. For a
hypothetical 10,000 accounts and 1,000 USD per unit outcome, the associated
expected loss ranges from 1,000,000 to 6,000,000 USD. Imputing one value
selects a restriction; it does not learn the missing population from the
observations alone. Public sensitivity tests can reproduce these endpoints,
while actual bank missingness needs further information or justified
restrictions \cite{exec24_sas_choice}.

The pooled covariance and transformation scale must also survive export.
The retained combining-inferences discussion belongs to the uncertainty
contract, not to an optional presentation layer. An agent should record
which imputation model, analysis model, assumptions, overrides and failed
fits contributed to a result before narrating a business recommendation
\cite{exec21_transport_sas_combining,exec21_transport_sas_mianalyze_intro,
exec24_sas_choice}.

\subsection{Terminal jobs and bounded polling are different guarantees}
The inspected sasctl generic job monitor treats completed and failed as
terminal states. Returning a terminal job therefore does not by itself
mean success. Its refresh limit counts requests and its loop sleeps half
a second per refresh; the selected request delegate separately accepts a
timeout argument whose default is None. This source inspection does not
prove that every service uses this helper, nor does it establish actual
server behavior \cite{exec23_workflow_delegates_sasctl}.

An authored timing model makes the required contract explicit. Sixty
polls with a half-second sleep and one-second request duration take
90 seconds; with sixty-second request duration they take 3,630 seconds.
A bound on poll count alone cannot bound wall time without a request-time
bound. A proposed worker therefore needs an end-to-end deadline, explicit
request limits and a typed unresolved state after local expiry. Cancellation
and recovery of the remote job require separately qualified semantics
\cite{exec24_sas_choice}.

The selected score-polling method has a different loop and return contract;
its timeout branch appears inside the loop, without an explicit return
after loop exit. The workflow class carries an experimental warning about
its partial implementation and historical documentation status. That
warning is preserved as a claim about the inspected source, not recast as
a verified statement about every present-day API. The request delegate's
token-refresh retry and its tests' mocked HTTP operations likewise require
public fault and recovery fixtures before native qualification
\cite{exec23_workflow_sasctl,exec23_workflow_delegates_sasctl}.

The historical Job Execution sample separates definitions, execution jobs,
states, errors, results and cleanup. Its result map is provider-specific.
Visual Investigator's BPMN case-management workflow is a different context;
it must not be substituted for a generic analytical workflow merely because
both use the word workflow. A CLI or MCP response must retain the server
profile, request identity, provider, terminal state and result schema
\cite{exec24_workflow_contracts_sas}.

\subsection{Pagination and complete costs matter to scientific results}
The inspected client pagination uses offsets and limits derived from links,
and its item count is documented as approximate. This does not establish
a snapshot-consistent analysis population. In an authored example, the
first page of the ordered list $(a,b,c,d)$ is $(a,b)$. Inserting $x$ at the
front before fetching offset two produces the second page $(b,c)$. Four
returned records contain three unique original accounts and omit $d$.
This conditional witness is not a native SAS failure; it specifies a
test for immutable extraction, stable identities and changing populations
\cite{exec23_workflow_delegates_sasctl,exec24_sas_choice}.

The reference route is retained for the named statistical and operational
contracts. Conditional integration requires a supported entitled procedure
or service plus public fixtures. Direct adoption as the complete causal
product is deferred. Open-client reuse or porting needs the actual selected
component's notices and dependencies; proprietary statistical implementations
cannot be ported on the basis of an open wrapper's license. Unsupported
causal, timing and success claims are rejected without rejecting every
bounded use of the components \cite{exec24_sas_choice}.

Scientific owners must specify targets, assumptions and loss; statistical
owners must specify event and missing-data laws; adapter owners must retain
model, time and extraction semantics; platform owners must establish exact
service support, rights and failure recovery. Hypothetical expenditure of
36,000 USD over 1,200 submitted, 900 completed and 600 qualified requests
gives 30, 40 and 60 USD respectively. Actual prices and economic value
remain unmeasured. Remaining treatment, panel, discovery and service scopes
need exact primary evidence; the present selection is not an exhaustive
inventory. Package execution and independent validation remain to be performed \cite{exec24_sas_choice}.


% SOURCE docs/release/2026-10-29/source/docs/writing-repair/2026-10-29/src/6d6b2e1570e0d761.tex CHAR 0
\section{Dara graph applications and the authority to change a causal model}
\label{sec:exec25-dara-reuse}
Dara provides a concrete application and graph-editing reference for the
proposed causal library. The inspected repository revision
\texttt{0ca24611ef35a9725f14a}\allowbreak\texttt{f6cda83718482665294}
is dated 30 September 2026. Its selected Python core, component and
application-generator packages, and JavaScript graph editor, declare version
1.29.11. Nineteen retained source files were authenticated against the complete
Git tree. The root and editor notices refer to Apache 2.0; dependencies,
assets, deployed services and commercial offerings remain separate rights and
qualification questions. This is an inspection of immutable source bytes,
not an installation or demonstration of a causal service.
\citep{exec25_dara_source}

The causaLens/Dara umbrella is resolved into distinct proposal decisions.
Dara's framework, graph viewer/editor and application CLI are retained as
references, with component-specific conditional integration. Chapter~\ref{scisyn:software} separately identifies the later advertised
causaLens service offering. Its commercial integration remains deferred until
a supported artifact, interface, applicable terms and substantive technical
evidence are selected; the open Dara source establishes none of those by itself.
Open framework code cannot establish proprietary product capabilities,
entitlement, pricing or operational acceptance. The proposal therefore neither
borrows commercial approval from the framework nor infers a commercial
absence from the selected open-source paths.~\citep{exec25_dara_choice}

A branch manager may want to revise a diagram after learning that an account
status was recorded after a servicing call. A useful interface can display the
graph and collect a proposed correction. Scientific acceptance additionally
requires the question version, source evidence, variable identity, observation
clock, edit authority and consequences for identification. A drag action or
confirmation flag is evidence of a user interaction; it does not prove that
the revised arrow represents a causal mechanism. This distinction is central
to using an application framework around the proposed causal engine.
\citep{exec25_dara_choice}

\subsection{Graph modes and conversions are scientific contracts}
The selected Python viewer accepts a causal-graph object, skeleton, dictionary
or variable containing graph data. Its local validator returns variables
unchanged and checks the declared input type for other inputs. The same class
changes a non-default editor mode to DEFAULT when a planar layout is selected,
with a warning. The source describes DEFAULT as a DAG-oriented viewer, PAG
mode as displaying additional endpoint types, and RESOLVER as supporting
edge acceptance. A proposed adapter must make such mode changes visible and
verify that the declared graph class and unresolved marks survive conversion.
A type check on a dictionary is not a full proof of causal-graph validity or
of the assumptions behind the displayed graph.~\citep{exec25_dara_source}

The inspected editor types distinguish directed, bidirected, undirected and
partially marked edges, including an internal backwards-directed form. These
symbols can carry materially different scientific meanings. A latent-variable
analysis that returns unresolved endpoint marks must not be displayed as one
fully directed intervention graph simply because a layout or editing mode
expects it. A round-trip oracle should compare identifiers, endpoint marks,
lag or time information where admitted, graph class, metadata provenance and
the target query, not just whether the picture contains the same number of
edges. Unsupported graph classes should produce an explicit refusal or a
clearly labeled visualization-only view.~\citep{exec25_dara_source,exec25_dara_choice}

One selected serialization composition illustrates the need for that oracle.
The edge serializer converts a backwards-directed representation to a directed
edge and exchanges its embedded source and destination. The enclosing graph
serializer again exchanges those embedded endpoints while placing the edge
under reversed dictionary keys. For a stipulated internal pair $(A,B)$, the
first exchange yields embedded $(B,A)$; the second yields $(A,B)$, while the
outer keys are $(B,A)$. This is an exact symbolic trace of the selected source
branches, not a claim that an installed frontend has produced a reachable
failure. It identifies a concrete consistency test: the endpoint identifiers
inside every serialized edge must agree with its enclosing source/target
keys under the declared convention. Both direct and backwards representations
must pass the same canonicalization test.~\citep{exec25_dara_source,exec25_dara_choice}

\subsection{Editing, acyclicity and evidence authority}
The default-mode edge-add path checks whether the proposed addition creates
a cycle. In the selected reverse-edge path, the original edge is removed from
a graph copy and the original selected tuple is supplied to the cycle helper;
the eventual reducer exchanges the edge direction. That argument sequence
requires a targeted test against the intended reversed graph. An authored
example starts from $A\to B$, $A\to C$, $C\to B$, which is acyclic. Reversing
$A\to B$ creates $B\to A\to C\to B$. Checking the original direction does
not certify the reversed one. The source trace and exact graph calculation
motivate a qualification fixture; no upstream code or native failure was
executed. The test must also respect mode-specific semantics: DAG acyclicity
is not an all-purpose validity criterion for every partial ancestral graph.
\citep{exec25_dara_source,exec25_dara_choice}

The selected reducer records edge acceptance in rendering metadata. The
Python/JavaScript wrapper updates its graph variable and invokes the update
action, while the editor coalesces outgoing updates with a trailing
150-millisecond debounce. These mechanisms support interaction, but a
coalesced callback does not by itself provide an authenticated record of every
scientific edit or a transaction approved by a causal backend. The proposed
integration should distinguish a local draft, a submitted change and an
accepted model version. A change prepared for graph version $g_1$ must be
reconciled or refused if the current graph is $g_2$; it must not silently
apply to a different model.~\citep{exec25_dara_source,exec25_dara_choice}

For the servicing example, an edit that removes a mediator from a proposed
adjustment set changes the analysis contract even if the screen still shows
the same nodes. The accepted record should identify the author or authorized
role, evidence, old and new graph digests, question version, restriction type,
reason and resulting identification status. An LLM can draft the explanation
or propose alternatives, while the service preserves unknown versus absent
edges and conditional versus identified effects. Rendering acceptance,
authority to edit and scientific justification are three separate fields.
\citep{exec25_dara_choice}

\subsection{Application CLI versus a causal service interface}
The retained metadata declares application-creation and application-runtime
CLI entry points. Their role is useful for packaging an interface; it does
not establish a CLI for causal estimands or an MCP service contract. The
selected component dependency on \texttt{cai-causal-graph>=0.3.6} also leaves
a concrete delegate and environment to qualify. Python and JavaScript
package versions, bundled assets, protocol schemas and dependencies must be
bound together before a deployed interface can be called reproducible.
A root notice does not settle the entire dependency or deployment inventory.
\citep{exec25_dara_source}

Conditional integration would place a versioned graph editor in front of the
library's question, identification and execution contracts. Requests would
carry graph class, identifiers and marks, question version, data/schema
digests, restrictions, edit provenance, authority, expected graph version and
a deadline. Responses would separate draft rendering from accepted graph
state, identified or conditional effects, refusal and execution status. A
port of selected editor code would still require rights, dependency and
semantic tests; it would not inherit a native causal guarantee from the
original UI. Direct adoption as a causal inference engine is not supported
by the selected application source.~\citep{exec25_dara_choice}

The proposed public qualification suite uses finite known graphs and generated
structural laws for scientific truth. It includes all admitted edge marks,
identifier permutations, backwards-edge serialization, the reversal-cycle
example, an attempted planar/PAG mode change, stale edits, rejected backend
changes, duplicate updates and late callbacks. A fixed graph-file workflow,
a permitted frozen Dara interface and the proposed adapter should receive
the same questions and information. The protocol must freeze schemas,
versions, graph conventions, scientific acceptance and complete denominators
before native measurement. Permitted public data can add realistic schema
and interaction stress without supplying unobserved causal truth. None of
these native trials is claimed as completed here.~\citep{exec25_dara_choice}

Scientific owners would approve graph restrictions, questions and
identification; interface owners would preserve rendering and serialization;
service owners would enforce version and authority checks; rights/platform
owners would select permitted artifacts; evaluators and maintainers would
control fixtures and regressions. These remain proposed responsibilities.
An authored USD 1,000 preparation cost, USD 1,000 of all attempts,
USD 2,500 scientific review and USD 500 maintenance allocation sum to
USD 5,000. For 100 attempts, 80 completions and 50 admissible outputs, costs
are USD 50, USD 62.50 and USD 100. These are accounting examples, not vendor
prices. The selected reference route remains conditional on rights,
semantic preservation and measured useful scientific outcomes.
\citep{exec25_dara_choice}


% SOURCE docs/release/2026-10-29/source/docs/writing-repair/2026-10-29/src/45c8c46a0c8ae1db.tex CHAR 0
\section{AgenaRisk: a public wrapper around a licensed causal-modeling workflow}
\label{sec:exec25-agena-reuse}
The public AgenaRisk API supplies a useful Java and command-line comparator,
but its identity must be stated precisely. The selected repository commit
\texttt{54115453b0c8d500fa45}\allowbreak\texttt{bdffaeaf728c825391d3},
dated 1 September 2026, is the most recent commit returned by the recorded
cutoff query. Its README announces that development moved to a private
repository, calls 1.042 the final public release, and describes a later binary
distribution. The same selected tree declares a 1.043-SNAPSHOT Maven project
and a core dependency named \texttt{master-SNAPSHOT}. These are distinct
identifiers. A frozen wrapper commit does not freeze the engine, native
libraries, license material or later private API. Sixteen original source
and rights files are authenticated to the complete selected Git tree.
\citep{exec25_agena_source}

The proposal retains this public wrapper as a reference, with conditional
integration of a separately permitted, fully identified commercial runtime.
The selected source is not a basis for porting proprietary core code or
asserting the capabilities of an uninspected private release. The root
Apache notice applies to published source, while the README separately
requires a valid API usage license. The retained Modeller agreement identifies
version 2.4 and a 16 July 2025 update date; section 2.5 describes restrictions
on embedding and network-accessible service use. Those recorded terms identify
questions for a deployment-specific agreement, rather than establishing
permission for the proposed MCP service. No agreement was accepted, licensed
runtime installed or vendor computation performed.~\citep{exec25_agena_source,exec25_agena_choice}

\subsection{Evidence propagation and intervention are different requests}
Consider a bank studying whether a servicing call prevents delinquency. The
observed fact that a customer received a call may reveal prior financial
distress. Forcing the call policy is a different operation. The selected
\texttt{DataSet} methods install hard evidence or weighted observations in a
scenario, and the model calculation delegates propagation to the core. The
inspected observation path checks model membership and state/type compatibility.
It does not itself replace the treatment mechanism or remove incoming causal
arrows. This describes the selected method, not a universal claim that another
vendor interface cannot represent interventions.~\citep{exec25_agena_source}

An exact authored example makes the distinction testable. Let latent distress
$U$ have probability one half, let $Y=U$ denote subsequent delinquency, and
let the treatment indicator satisfy $P(A=1\mid U=1)=4/5$ and
$P(A=1\mid U=0)=1/5$. Here the call has no effect on $Y$. Conditioning gives
$E[Y\mid A=1]=4/5$ and $E[Y\mid A=0]=1/5$, an observed contrast of $3/5$.
Replacing the assignment mechanism by either fixed action leaves
$E[Y\mid\operatorname{do}(A=a)]=1/2$ for both actions, giving zero effect.
The calculation enumerates the finite law with exact fractions; it is not
an AgenaRisk output. A proposed adapter must distinguish evidence, intervention
and counterfactual requests and refuse a causal interpretation when only an
observational calculation has been justified.~\citep{exec25_agena_choice}

The selected calculation flags also deserve explicit exposure. They request
ancestor-network calculation, retention of tails and zero-mass regions, or
the log probability of evidence. The wrapper maps them to corresponding core
flags. Retaining numerical tails is not a proof of treatment positivity, and
the probability of evidence is not the value of a policy. The model, parameter
tables, variable states, observation schema, selected flags and target query
must travel together in the analysis record. Numerical agreement should be
checked against an independent finite-network oracle before a result is
interpreted as scientific support.~\citep{exec25_agena_source,exec25_agena_choice}

\subsection{A restartable CLI requires more than a scenario identifier}
The retained calculator accepts a model file, dataset file and output file.
Its optional cache reads previous JSON result objects, collects their dataset
identifiers and skips matching identifiers on a later run. The inspected
selection predicate does not compare model, data, engine or query digests.
After each successful calculation it rewrites the accumulated JSON array with
\texttt{TRUNCATE\_EXISTING}. That path establishes useful incremental output,
but does not itself establish an atomic replacement protocol or content-bound
cache validity. These are static source findings, not observed corruption or
stale results from an installed process.~\citep{exec25_agena_source}

Suppose branch scenario \texttt{branch\_17} is retained while its model changes
from $m_1$ to $m_2$ and data from $d_1$ to $d_2$. Identifier equality still holds,
even if stipulated correct probabilities change from $1/5$ to $4/5$. An
identifier-only skip would retain a value differing by $3/5$ in this authored
witness. A proposed cache key must bind the canonical model, parameters,
data/schema, wrapper and engine versions, flags, target, scientific restrictions
and output schema. Repeated requests also need a distinct attempt history;
idempotent reuse must not erase failures or costs. Fault trials should interrupt
the write at controlled points and verify that a complete previous or new
artifact survives, or that the incomplete state is explicitly rejected.
\citep{exec25_agena_choice}

The wrapper separates inconsistent evidence, resource failures and general
calculation exceptions after delegated propagation. Its selected method
temporarily modifies core message settings and captures output, restoring
settings in a finally block. This motivates concurrency and process-isolation
fixtures; it does not prove a race in a deployed service. The CLI normally
exits on a dataset exception or error, using different exit codes, and its
launcher separately reports setup errors. A service must map these paths to
typed outcomes and preserve every requested scenario, including one that
produces no result object. No overall deadline, safe retry policy or MCP
contract follows merely from having a command-line entry point.
\citep{exec25_agena_source,exec25_agena_choice}

\subsection{Learning delegates, public qualification and maintenance}
The structure-learning launcher reads a configuration file and delegates to
\texttt{StructureLearner}; the Maven metadata names a separate learning
artifact. The launcher exposes thread configuration, but does not by itself
establish the learning algorithm, search completeness, graph equivalence
class, hidden-confounding assumptions or validity under dependent observations.
Those require the matching delegate and configuration to be inspected. The
proposed library must retain graph uncertainty and connect any discovered
object to the assumptions needed for the specific intervention query.
\citep{exec25_agena_source,exec25_agena_choice}

The build metadata selects a mutable core snapshot and downloads native or
product resources during the clean phase. Consequently even a pinned public
wrapper is insufficient for reproducible installation. Qualification should
first select permitted immutable artifacts and a deployment profile, then
freeze all dependencies and rights, schemas, numeric settings and budgets.
The README's distinction between ordinary and enterprise/headless activation
is another reason to specify the target deployment before pricing or measuring
an integration. Current private-release support and actual commercial prices
remain unverified.~\citep{exec25_agena_source}

The public qualification plan compares a finite-network reference calculation,
a permitted pinned runtime and the proposed adapter under identical questions
and information. Fixtures cover evidence versus intervention, impossible
evidence, state renaming, numerical tails, stale same-ID caches, duplicate
requests, interrupted writes, resource refusal, process cancellation, concurrent
calls and missing result objects. Generated structural laws provide known
causal truth; permitted public datasets provide realistic schema and workflow
stress without supplying unobserved truth. Learning trials separately require
the exact delegate, constraints, data split and graph/effect scoring rules.
These trials are proposed, not executed.~\citep{exec25_agena_choice}

Scientific reviewers would own model assumptions, estimands and evidence
interpretation; interface and platform owners would own schemas, process
isolation and recovery; rights owners would obtain the selected entitlement;
maintainers would manage the frozen-public/private-product transition. These
are proposed responsibilities. An authored USD 1,200 preparation cost,
USD 1,800 for all attempts, USD 4,000 scientific review and USD 1,000 maintenance
allocation total USD 8,000. With 200 attempts, 160 completions and 100 admissible
results, the costs are USD 40, USD 50 and USD 80 respectively. None is a vendor
quote. The reference choice closes a bounded comparison decision while leaving
the complete runtime, rights, scientific qualification and economic case open.
\citep{exec25_agena_choice}

\section{The result must survive the complete workflow}

The commercial comparisons support a specific recommendation: select an
operation only after defining the information and target it preserves. An
adjustment set, a sharp bound, a model-conditional counterfactual, a cohort
effect and a policy value are useful but distinct outputs. The exact examples
show why observed fit, predictive ranking, population intervention agreement
or a successful numerical call can still leave a banking action unresolved.
The same distinction applies to uncertainty: outcome variance, posterior
spread, sampling error and ambiguity across models cannot share an unlabeled
interval field.

The product studies carry that argument into actual boundaries. State order,
graph marks, parameter status, event clocks, selected populations, licenses,
asynchronous failures and cache identity can all change the interpretation
or cost of an otherwise correct calculation. Their implementation-specific
witnesses remain reasons to test a narrow adapter, rather than grounds for a
product-wide ranking. Existing engines supply substantial capabilities; the
proposed common layer earns its place only if the complete comparison shows
that it preserves meaning and improves the accepted workflow. The next
chapter develops the corresponding obligations when computation crosses a
language boundary.

